\documentclass[11pt]{article}
\usepackage{orcidlink}
\usepackage[margin=0.75in]{geometry}
\usepackage{amsmath,amssymb,amsfonts,bm,mathtools}
\usepackage{mathrsfs}
\usepackage{booktabs,longtable,array,multirow}
\usepackage{xcolor}
\usepackage{graphicx}
\usepackage{hyperref}
\usepackage[T1]{fontenc}
\usepackage{lmodern}
\usepackage[utf8]{inputenc}
\usepackage{enumitem}
\usepackage{amsmath}
\usepackage[most]{tcolorbox}
\usepackage{soul}
\usepackage{needspace}
\usepackage{float}
\usepackage[font=small,labelfont=bf]{caption}


\setlength{\parindent}{0pt}
\setlength{\parskip}{0.55em}
\renewcommand{\arraystretch}{1.22}

\newcommand{\PR}{\mathrm{PR}}
\newcommand{\eff}{\mathrm{eff}}
\newcommand{\Kerr}{\mathrm{Kerr}}
\newcommand{\GR}{\mathrm{GR}}
\newcommand{\obs}{\mathrm{obs}}
\newcommand{\src}{\mathrm{src}}
\newcommand{\vac}{\mathrm{vac}}
\newcommand{\homPR}{\mathrm{hom}}
\newcommand{\TT}{\mathrm{TT}}
\newcommand{\Teuk}{\mathrm{Teuk}}
\newcommand{\dd}{\mathrm{d}}
\newcommand{\Tr}{\mathrm{Tr}}
\newcommand{\diag}{\mathrm{diag}}
\renewcommand{\arraystretch}{1.4}

\title{Supplementary Material: \\ 
\Large Extended Derivations and Automated Verification \\ 
\textit{The Foundation of Projection Relativity}}

\author{Michael Stanislaus Oshetski \orcidlink{0009-0007-3623-7586} \\[-6pt]
\small \email{michael.oshetski@micatu.com}\\[-3pt]
\small Founder and Chief Technology Officer \\[-6pt]
\small Advanced Photonics Laboratory (APL)\\[-6pt]
\small Micatu, Inc. \\[-6pt]
\small Horseheads, New York, USA \\[-6pt]
\small ORCID: 0009-0007-3623-7586}
\date{May 2026}

\begin{document}
\maketitle

\section{Introduction to the Supplemental Material}

This supplementary document aggregates the foundational equations of Projection Relativity, integrating them with the complete automated verification logs. It provides the exhaustive, step-by-step mathematical derivations that support the core framework, serving both as a detailed analytical expansion of the main manuscript and a formal repository for the computational test harness results.

The formalisms and algorithmic outputs presented herein detail the internal metric tensor, spectral operators, modal expansions, projection maps, compact boundary closure conditions, and the derived-status matrix. Crucially, the appended symbolic regression transcripts and numerical testing suites establish rigorous computational validation for the spectral gaps, topological invariants, and macroscopic boundary conditions derived in the text.

\renewcommand{\contentsname}{Table of Contents} % Optional: Renames "Table of Contents" to "List of Sections"
\tableofcontents
\clearpage % Pushes the actual math content to the next page


\section{Master Equation Set}
\label{app:how-to-use}

This section collects the core equations of Projection Relativity in one place. It is a reference map for the main derivations, not a separate derivation chain. The equations below summarize the internal metric, spectral operators, mode expansions, projection maps, compact boundary closure, and derived-status table used throughout the paper.

The organizing structre is
\begin{equation}
\boxed{
\Psi
\rightarrow
O_{\rm int}
\rightarrow
\{U_{n,m},\Lambda_{n,m}\}
\rightarrow
\{P_g,P_{\rm disp},P_A,P_X,P_\theta\}
\rightarrow
\text{observable projection sectors}.
}
\end{equation}
The detailed sector derivations are given in the main text and in the later appendices. This reference sheet keeps the master structure visible.

\subsection{Master Equation Matrix Expansion}
\label{app:matrix-expansion}

This subsection gives the full Projection Relativity matrix expansion used for visualization, bookkeeping, and linear analysis of the framework constructed in Section~\ref{sec:foundational_mathematical_structures}.

\subsubsection{Internal Metric, Measure, and Hilbert Space}
\label{app:internal-metric}

The internal manifold is
\begin{equation}
\boxed{
\mathcal M_{\rm int}
=
\mathbb R_w\times S^1_\theta,
\qquad
\xi^A=(w,\theta).
}
\end{equation}
The minimal internal metric is
\begin{equation}
\boxed{
G_{AB}
=
\begin{pmatrix}
1 & 0 \\
0 & R_A^2
\end{pmatrix}.
}
\end{equation}
The inverse metric is
\begin{equation}
\boxed{
G^{AB}
=
\begin{pmatrix}
1 & 0 \\
0 & R_A^{-2}
\end{pmatrix}.
}
\end{equation}
The determinant is
\begin{equation}
\boxed{
\det(G_{AB})=R_A^2.
}
\end{equation}
The invariant internal measure is
\begin{equation}
\boxed{
d\mu_{\rm int}
=
\sqrt{\det(G_{AB})}\,dw\,d\theta
=
R_A\,dw\,d\theta.
}
\end{equation}
The projection Hilbert space is
\begin{equation}
\boxed{
\mathcal H_P
=
L^2(\mathcal M_{\rm int},d\mu_{\rm int}).
}
\end{equation}
The corresponding inner product is
\begin{equation}
\boxed{
\langle f|g\rangle_P
=
\int_{\mathcal M_{\rm int}}
d\mu_{\rm int}\,
f^*(\xi)g(\xi).
}
\end{equation}

\subsubsection{Laplace--Beltrami Operator}
\label{app:laplace-beltrami}

The scalar Laplace--Beltrami operator on the internal manifold is
\begin{equation}
\boxed{
\Delta_G
=
\frac{1}{\sqrt{\det(G)}}
\partial_A
\left(
\sqrt{\det(G)}
G^{AB}
\partial_B
\right).
}
\end{equation}
For stationary compact radius,
\begin{equation}
\boxed{
\partial_wR_A=0,
\qquad
\partial_\theta R_A=0,
}
\end{equation}
this reduces to
\begin{equation}
\boxed{
\Delta_G
=
\partial_w^2
+
R_A^{-2}\partial_\theta^2.
}
\end{equation}

\subsubsection{Radial Spectral Operator}
\label{app:radial-operator}

The radial spectral operator is
\begin{equation}
\boxed{
O_X
=
-\frac{d^2}{dw^2}
+
V_X(w).
}
\end{equation}
The positive quartic family is
\begin{equation}
\boxed{
V_X(w)
=
\Lambda_X^2
\left(
1+a_2w^2+a_4w^4
\right).
}
\end{equation}
The Projection Relativity trace-selected radial branch is
\begin{equation}
\boxed{
V_{X,\star}(w)
=
\Lambda_X^2
\left(
1+w^2+\frac{3}{4}w^4
\right).
}
\end{equation}
In the reference normalization \(\Lambda_X^2=1\), this is
\begin{equation}
\boxed{
V_{X,\star}(w)
=
1+w^2+\frac{3}{4}w^4.
}
\end{equation}
The radial eigenvalue equation is
\begin{equation}
\boxed{
O_Xu_n
=
\lambda_nu_n.
}
\end{equation}

\subsubsection{Compact Phase Operator}
\label{app:compact-phase-operator}

The compact phase operator is
\begin{equation}
\boxed{
O_\theta
=
-\frac{1}{R_A^2}
\frac{d^2}{d\theta^2}.
}
\end{equation}
The normalized compact eigenmodes are
\begin{equation}
\boxed{
v_m(\theta)
=
\frac{1}{\sqrt{2\pi R_A}}
e^{im\theta},
\qquad
m\in\mathbb Z.
}
\end{equation}
The compact eigenvalues are
\begin{equation}
\boxed{
\lambda_m^{(\theta)}
=
\frac{m^2}{R_A^2}.
}
\end{equation}

\subsubsection{Full Internal Spectral Operator}
\label{app:internal-spectral-operator}

The full internal operator is separable:
\begin{equation}
\boxed{
O_{\rm int}
=
O_X+O_\theta.
}
\end{equation}
In expanded form,
\begin{equation}
\boxed{
O_{\rm int}
=
-\frac{d^2}{dw^2}
-
\frac{1}{R_A^2}
\frac{d^2}{d\theta^2}
+
V_X(w).
}
\end{equation}
The internal eigenmodes factorize as
\begin{equation}
\boxed{
U_{n,m}(w,\theta)
=
u_n(w)
\frac{e^{im\theta}}{\sqrt{2\pi R_A}}.
}
\end{equation}
The corresponding eigenvalues are
\begin{equation}
\boxed{
\Lambda_{n,m}
=
\lambda_n
+
\frac{m^2}{R_A^2}.
}
\end{equation}
A schematic matrix representation is
\begin{equation}
\boxed{
O_{\rm int}
=
\begin{pmatrix}
\Lambda_{0,0} & 0 & 0 & \cdots \\
0 & \Lambda_{0,1} & 0 & \cdots \\
0 & 0 & \Lambda_{1,0} & \cdots \\
\vdots & \vdots & \vdots & \ddots
\end{pmatrix}.
}
\end{equation}

\subsubsection{Master Field Expansion}
\label{app:master-field-expansion}

The master field expands in the internal spectral basis as
\begin{equation}
\boxed{
\Psi(x,w,\theta)
=
\sum_{n,m}
c_{n,m}(x)
U_{n,m}(w,\theta).
}
\end{equation}
Equivalently, in schematic coefficient-vector form,
\begin{equation}
\boxed{
\Psi
=
\begin{pmatrix}
\vdots \\
c_{0,0}(x)U_{0,0} \\
c_{0,1}(x)U_{0,1} \\
c_{1,0}(x)U_{1,0} \\
c_{1,1}(x)U_{1,1} \\
\vdots
\end{pmatrix}.
}
\end{equation}

\subsubsection{Master Action}
\label{app:master-action}

Using the mostly-plus convention used in the main text, the master action is
\begin{equation}
\boxed{
S_{\rm PR}
=
\int d^4x\,d\mu_{\rm int}
\sqrt{-g_{\rm eff}}
\left[
-\frac12
g_{\rm eff}^{\mu\nu}
D_\mu\Psi^\dagger
D_\nu\Psi
-
\frac12
G^{AB}
D_A\Psi^\dagger
D_B\Psi
-
V_{\rm PR}
+
\mathcal L_{\rm src}
\right].
}
\end{equation}
The internal kinetic structure is
\begin{equation}
\boxed{
G^{AB}
D_A\Psi^\dagger D_B\Psi
=
D_w\Psi^\dagger D_w\Psi
+
R_A^{-2}
D_\theta\Psi^\dagger D_\theta\Psi.
}
\end{equation}

\subsubsection{Euler--Lagrange Spectral Equation}
\label{app:euler}

In the quadratic spectral regime, the master equation is
\begin{equation}
\boxed{
(\square+O_{\rm int})\Psi=0.
}
\end{equation}
The internal spectral equation is
\begin{equation}
\boxed{
O_{\rm int}U_{n,m}
=
\Lambda_{n,m}U_{n,m}.
}
\end{equation}
Substituting the mode expansion gives
\begin{equation}
\boxed{
\sum_{n,m}
\left[
(\square c_{n,m})
+
\Lambda_{n,m}c_{n,m}
\right]
U_{n,m}
=
0.
}
\end{equation}
Projecting onto each orthonormal internal mode gives the observable spacetime equation
\begin{equation}
\boxed{
\left(
\square
-
\Lambda_{n,m}^{\rm phys}
\right)c_{n,m}(x)
=
0.
}
\end{equation}

\subsubsection{Projection-Sector Dictionary}
\label{app:projection-dictionary}

The principal projection-sector maps are
\begin{align}
P_g[\Psi]
&\rightarrow
g_{\mu\nu}^{\rm eff},
\\
P_{\rm disp}[\Psi]
&\rightarrow
\Phi_{\rm disp}(x)
\rightarrow
A(x)
\rightarrow
A_{\rm phys}(x),
\\
P_A[\Psi]
&\rightarrow
\theta(x)
\rightarrow
\mathcal A_\mu(x),
\\
P_X[\Psi]
&\rightarrow
\Sigma_X,
\\
P_\theta[\Psi]
&\rightarrow
H_{\rm eff}(t).
\end{align}

\subsubsection{Spectral Gap Structure}
\label{app:spectral-gap}

The first radial spectral gap is
\begin{equation}
\boxed{
\mu_{\min}^2
=
\lambda_1-\lambda_0.
}
\end{equation}
In the reference radial branch,
\begin{equation}
\boxed{
\lambda_0
=
2.322863529580,
}
\end{equation}
\begin{equation}
\boxed{
\lambda_1
=
5.375830272676,
}
\end{equation}
and
\begin{equation}
\boxed{
\mu_{\min}^2
=
3.052966743096.
}
\end{equation}
The corresponding positive gap scale is
\begin{equation}
\boxed{
\mu_{\min}
=
1.747274089288.
}
\end{equation}
The fourth-power suppression scale is
\begin{equation}
\boxed{
\mu_{\min}^4
=
9.320605934450.
}
\end{equation}

\subsubsection{Compact Boundary-Orientation Fine-Structure Closure}
\label{app:compact-boundary-closure}

This subsection gives the reference form of the boundary-resolved compact phase closure used in Section~\ref{sec:compact_phase_closure_alpha}. The compact boundary cofactor is not a calibrated correction. It is the terminal-word cofactor of the finite-rank compact return map.

\subsubsection{Trace alphabet}
\label{app:trace-alphabet}

The compact and spatial trace fractions are
\begin{equation}
\boxed{
T_A=\frac14,
\qquad
T_X=\frac34.
}
\end{equation}
Compact returns are written as words in the symbol \(A\), weighted by powers of \(T_A\). Finite-rank boundary crossings are written as words containing one symbol \(X\), weighted by \(T_X\). The finite-rank slice is
\begin{equation}
\boxed{
\Pi_{13}
=
|\phi_1\rangle\langle\phi_1|
+
|\phi_3\rangle\langle\phi_3|.
}
\end{equation}

\subsubsection{Interior compact recurrence.}
\label{app:compact-recurrence}

The compact interior has two primitive return loops, \(A^3\) and \(A^4\). The corresponding transfer matrix is
\begin{equation}
\boxed{
M_A
=
\begin{pmatrix}
T_A^3&T_A^4\\
1&0
\end{pmatrix}.
}
\end{equation}
The Fredholm determinant is
\begin{equation}
\boxed{
D_A
=
\det(I-M_A)
=
1-T_A^3-T_A^4.
}
\end{equation}
For \(T_A=1/4\),
\begin{equation}
\boxed{
D_A
=
1-\frac1{64}-\frac1{256}
=
\frac{251}{256}.
}
\end{equation}
The inverse determinant \(D_A^{-1}\) resums recurrent compact interior paths. Terminal boundary words cannot include words already generated by this recurrent interior inverse.

\subsubsection{Admissible terminal words.}
\label{app:terminal-words}

The finite-rank boundary admits the signed terminal words
\begin{equation}
\boxed{
\operatorname{Adm}_{\Pi_{13}}
=
\{\epsilon,\ A^3,\ A^4,\ XA^5,\ -XA^8\}.
}
\end{equation}
Their weights are
\begin{equation}
\boxed{
1,
\qquad
T_A^3,
\qquad
T_A^4,
\qquad
T_XT_A^5,
\qquad
-T_XT_A^8.
}
\end{equation}
The terms have the following roles:
\begin{itemize}
\item \(\epsilon\) is direct boundary passage.
\item \(A^3\) is the primitive spatial-trace compact return.
\item \(A^4\) is the primitive full-trace compact return.
\item \(XA^5=XA^{2+3}\) is the first compact leave--return pair \(A^2\) followed by the spatial primitive return \(A^3\), weighted by \(X\).
\item \(-XA^8=-XA^{4+4}\) is the finite-rank exclusion of double full-trace recirculation outside the \(\Pi_{13}\) slice.
\end{itemize}

\subsubsection{Uniqueness of the terminal cofactor.}
\label{app:terminal-cofactor-uniqueness}

The admissible set is unique within the tested finite-rank compact return class under compact closure, finite-rank terminality, parity/orientation, and inclusion--exclusion:
\begin{itemize}
\item \(X\), \(XA\), and \(XA^2\) are not compactly closed returns.
\item \(XA^3\) and \(XA^4\) are not new boundary leakage terms; their compact parts are already primitive cofactor residues.
\item \(XA^6=XA^{2+4}\) is an orientation-neutral full-trace boundary pairing and is absorbed into the recurrent Fredholm denominator.
\item \(XA^7=XA^{3+4}\) factors into primitive interior returns and is determinant-resummed.
\item \(XA^9, XA^{10},\ldots\) factor through recurrent compact interior paths or higher boundary recirculations and belong to \(D_A^{-1}\), not to the terminal cofactor.
\end{itemize}

\subsubsection{Terminal cofactor.}
\label{app:terminal-cofactor}

The terminal cofactor is
\begin{equation}
\boxed{
N_{\rm bc}
=
1+T_A^3+T_A^4+T_X(T_A^5-T_A^8).
}
\end{equation}
For \(T_A=1/4\) and \(T_X=3/4\),
\begin{equation}
\boxed{
N_{\rm bc}
=
1+\frac1{64}+\frac1{256}
+
\frac34\left(\frac1{1024}-\frac1{65536}\right)
=
\frac{267453}{262144}.
}
\end{equation}

\subsubsection{Boundary-resolved compact return.}
\label{app:boundary-return}

The compact boundary return is the cofactor divided by the Fredholm determinant:
\begin{equation}
\boxed{
R_{\rm bc}
=
\frac{N_{\rm bc}}{D_A}
=
\frac{
1+T_A^3+T_A^4+T_X(T_A^5-T_A^8)
}{
1-T_A^3-T_A^4
}.
}
\end{equation}
The compact boundary constant entering the finite-rank spatial projection is
\begin{equation}
\boxed{
c_{\rm bc}
=
T_X
\left[
1+T_A^2
-
T_A^6R_{\rm bc}
\right].
}
\end{equation}
Equivalently,
\begin{equation}
\boxed{
c_{\rm bc}
=
T_X
\left[
1+T_A^2
-
T_A^6
\frac{
1+T_A^3+T_A^4+T_X(T_A^5-T_A^8)
}{
1-T_A^3-T_A^4
}
\right].
}
\end{equation}
Substitution gives
\begin{equation}
\boxed{
c_{\rm bc}
=
\frac{3354902985}{4211081216}
=
0.7966844648478991.
}
\end{equation}

\subsubsection{Finite-rank projection and leakage.}
\label{app:finite-rank-leakage}

The boundary-resolved projection operator is
\begin{equation}
\boxed{
P_{\rm geom}^{\rm bc}
=
\Pi_{13}K_X^2
\left(
c_{\rm bc}+w^2+\frac34w^4
\right)
K_X^2\Pi_{13}.
}
\end{equation}
The normalized stiffness operator is
\begin{equation}
\boxed{
K_X
=
\frac{O_X-\lambda_0}{\lambda_1-\lambda_0}.
}
\end{equation}
The finite-rank leakage probabilities are
\begin{equation}
\boxed{
p_1=7.6528903366\times10^{-4},
\qquad
p_3=1-p_1.
}
\end{equation}
The compact closure requires the finite-rank radial datum
\begin{equation}
\boxed{
\lambda_3
=
13.23388439508096.
}
\end{equation}
With
\begin{equation}
\boxed{
q_1=\lambda_1-\lambda_0,
\qquad
q_3=\lambda_3-\lambda_0,
}
\end{equation}
the effective compact stiffness is
\begin{equation}
\boxed{
q_{\rm bc}
=
R_{A,\star}^{-2}
=
p_1q_1+(1-p_1)q_3.
}
\end{equation}
Numerically,
\begin{equation}
\boxed{
q_{\rm bc}
=
10.905007182855176.
}
\end{equation}
The boundary-resolved compact normalization is
\begin{equation}
\boxed{
\alpha_{{\rm PR},bc}^{-1}
=
4\pi q_{\rm bc}
=
137.036361812007.
}
\end{equation}
This value is an output of the boundary-resolved compact projection chain. It is not used to choose \(c_{\rm bc}\), \(p_1\), \(R_A\), or \(C_A\).

\subsubsection{Full Master Chain}
\label{app:full-master-chain}

The complete reference chain is
\begin{equation}
\boxed{
\Psi
\rightarrow
O_{\rm int}
\rightarrow
\{U_{n,m},\Lambda_{n,m}\}
\rightarrow
\{P_g,P_{\rm disp},P_A,P_X,P_\theta\}
\rightarrow
\{g_{\mu\nu}^{\rm eff},A_{\rm phys},\mathcal A_\mu,\Sigma_X,H_{\rm eff}\}.
}
\end{equation}

\subsection{Status of Assumptions and Derived Quantities}
\label{app:status}

The assumptions, constraints, projection maps, and derived quantities introduced in Sections~\ref{sec:master_projection_field_foundational_postulates} and~\ref{sec:foundational_mathematical_structures} are summarized below. Quantities labeled ``Postulated'' are part of the minimal PR foundation. Quantities labeled ``Derived'' follow from the internal geometry, projection Hilbert space, projection trace, compact boundary map, or spectral operator structure. Quantities labeled ``Numerically derived'' or ``Numerically evaluated derived output'' are obtained by evaluating the corresponding derived operator without fitting the target observable.

\providecommand{\tblcell}[1]{\begin{minipage}[t]{\linewidth}\raggedright #1\end{minipage}}

\begingroup
\small
\setlength{\tabcolsep}{3pt}
\renewcommand{\arraystretch}{1.28}

\begin{longtable}{@{}>{\raggedright\arraybackslash}p{0.20\textwidth}>{\raggedright\arraybackslash}p{0.17\textwidth}>{\raggedright\arraybackslash}p{0.57\textwidth}@{}}
\caption{Status of assumptions and derived quantities in the PR master-equation architecture.}
\label{tab:status-assumptions-derived}
\\
\toprule
\textbf{Quantity} & \textbf{Status} & \textbf{PR-native origin} \\
\midrule
\endfirsthead

\toprule
\textbf{Quantity} & \textbf{Status} & \textbf{PR-native origin} \\
\midrule
\endhead

\midrule
\multicolumn{3}{r}{\emph{continued on next page}}\\
\midrule
\endfoot

\bottomrule
\endlastfoot

\(\Psi(x,\xi)\) &
Postulated &
\tblcell{Master projection field defined over observable spacetime and the internal spectral manifold.} \\

\(\xi^A=(w,\theta)\) &
Postulated coordinates &
\tblcell{Internal coordinates consisting of radial spectral coordinate \(w\) and compact phase coordinate \(\theta\).} \\

\(\mathcal M_{\rm int}\) &
Postulated topology &
\tblcell{Minimal internal projection manifold supplying radial stiffness and compact phase winding:
\[
\mathcal M_{\rm int}=\mathbb R_w\times S^1_\theta .
\]} \\

\(\theta\sim\theta+2\pi n\) &
Topological constraint &
\tblcell{Compact phase identification defining the \(S^1_\theta\) projection fiber, with \(n\in\mathbb Z\).} \\

\(G_{AB}\) &
Derived geometry &
\tblcell{Internal metric tensor obtained from the separable cylindrical line element:
\[
G_{AB}=\mathrm{diag}(1,R_A^2),
\qquad
ds_{\rm int}^2=dw^2+R_A^2d\theta^2 .
\]} \\

\(G^{AB}\) &
Derived geometry &
\tblcell{Inverse internal metric used in the internal Laplace--Beltrami operator:
\[
G^{AB}=\mathrm{diag}(1,R_A^{-2}).
\]} \\

\(d\mu_{\rm int}\) &
Derived geometry &
\tblcell{Invariant internal measure derived from \(\sqrt{\det(G_{AB})}\):
\[
d\mu_{\rm int}=R_A\,dw\,d\theta .
\]} \\

\(\mathcal H_P\) &
Derived Hilbert structure &
\tblcell{Projection Hilbert space used for internal normalization, overlap, and spectral expansion:
\[
\mathcal H_P=L^2(\mathcal M_{\rm int},d\mu_{\rm int}).
\]} \\

\(\langle f|g\rangle_P\) &
Derived Hilbert structure &
\tblcell{Internal inner product fixed by the invariant measure:
\[
\langle f|g\rangle_P
=
\int_{\mathcal M_{\rm int}}d\mu_{\rm int}\,f^*(\xi)g(\xi).
\]} \\

\(O_X\) &
Postulated radial operator &
\tblcell{Minimal radial spectral operator defining internal stiffness and confinement:
\[
O_X=-\frac{d^2}{dw^2}+V_X(w).
\]} \\

\(V_{X,\star}(w)\) &
Derived projection-trace branch &
\tblcell{Projection-trace radial branch:
\[
V_{X,\star}(w)
=
\Lambda_X^2\left(1+w^2+\frac34w^4\right).
\]
The quadratic coefficient \(a_2=1\) follows from unit local radial stiffness; the quartic coefficient \(a_4=3/4\) follows from the \(3+1\) projection trace. They are derived branch coefficients, not selected fit parameters.} \\

\(a_2=1\) &
Derived coefficient &
\tblcell{Unit local radial stiffness of the projection-trace radial branch.} \\

\(a_4=3/4\) &
Derived coefficient &
\tblcell{Projection-trace coefficient from the observable spatial trace relative to the full \(3+1\) trace.} \\

\(O_\theta\) &
Derived compact operator &
\tblcell{Compact phase operator obtained from the \(S^1_\theta\) Laplace--Beltrami sector:
\[
O_\theta=-\frac{1}{R_A^2}\frac{d^2}{d\theta^2}.
\]} \\

\(O_{\rm int}\) &
Derived internal operator &
\tblcell{Separable internal spectral operator combining radial stiffness and compact phase winding:
\[
O_{\rm int}=O_X+O_\theta .
\]} \\

\(u_n(w)\) &
Derived spectral mode &
\tblcell{Radial eigenmodes satisfying \(O_Xu_n=\lambda_nu_n\).} \\

\(v_m(\theta)\) &
Derived analytic mode &
\tblcell{Compact phase eigenmodes satisfying periodicity on \(S^1_\theta\):
\[
v_m(\theta)=\frac{1}{\sqrt{2\pi R_A}}e^{im\theta}.
\]} \\

\(\lambda_m^{(\theta)}\) &
Derived analytic eigenvalue &
\tblcell{Compact phase winding eigenvalues obtained from \(O_\theta v_m=\lambda_m^{(\theta)}v_m\):
\[
\lambda_m^{(\theta)}=\frac{m^2}{R_A^2}.
\]} \\

\(U_{n,m}\) &
Derived spectral basis &
\tblcell{Full separable internal spectral basis built from radial and compact eigenmodes:
\[
U_{n,m}(w,\theta)=u_n(w)v_m(\theta).
\]} \\

\(\Lambda_{n,m}\) &
Derived spectral eigenvalue &
\tblcell{Full internal eigenvalue obtained from separability of \(O_{\rm int}\):
\[
\Lambda_{n,m}=\lambda_n+\frac{m^2}{R_A^2}.
\]} \\

\(\lambda_0\) &
Numerically derived &
\tblcell{Ground-state radial eigenvalue from the stationary radial spectral operator:
\[
\lambda_0=2.322863529580.
\]} \\

\(\lambda_1\) &
Numerically derived &
\tblcell{First excited radial eigenvalue from the stationary radial spectral operator:
\[
\lambda_1=5.375830272676.
\]} \\

\(\lambda_3\) &
Numerically evaluated derived output &
\tblcell{Finite-rank radial spectral datum required by the compact boundary-resolved closure:
\[
\lambda_3=13.23388439508096.
\]
Together with \(p_1\), \(\lambda_0\), and \(\lambda_1\), this gives
\[
q_{\rm bc}
=
p_1(\lambda_1-\lambda_0)
+
(1-p_1)(\lambda_3-\lambda_0)
=
10.905007182855176.
\]} \\

\(\mu_{\min}^2\) &
Numerically derived &
\tblcell{First radial spectral gap; universal PR stiffness scale for finite-core regularization and low-energy suppression:
\[
\mu_{\min}^2=\lambda_1-\lambda_0=3.052966743096.
\]} \\

\(R_{\max}\) &
Derived hierarchy scale &
\tblcell{Maximum projection-curvature scale induced by the radial spectral stiffness scale \(\mu_{\min}^2\).} \\

\(r_c\) &
Derived hierarchy scale &
\tblcell{Finite-core saturation radius obtained from the hierarchy \(\mu_{\min}^2\rightarrow R_{\max}\rightarrow r_c\):
\[
r_c=\left(\frac{G_NM}{c^2R_{\max}}\right)^{1/3}.
\]} \\

\(\{P_i\}\) &
Defined projection maps &
\tblcell{Observable sector projections acting on the single master field:
\[
\{P_i\}=\{P_g,P_{\rm disp},P_A,P_X,P_\theta\}.
\]} \\

\(P_g[\Psi]\) &
Derived projection sector &
\tblcell{Gravitational stiffness projection generating the effective spacetime metric:
\[
P_g[\Psi]\rightarrow g_{\mu\nu}^{\rm eff}.
\]} \\

\(P_{\rm disp}[\Psi]\) &
Derived projection sector &
\tblcell{Complex displacement projection amplitude extracted from the master field:
\[
P_{\rm disp}[\Psi]\rightarrow\Phi_{\rm disp}(x).
\]} \\

\(\Phi_{\rm disp}(x)\) &
Derived projection structure &
\tblcell{Amplitude--phase decomposition of the displacement-projected field:
\[
\Phi_{\rm disp}(x)=A(x)e^{i\theta(x)}.
\]} \\

\(A(x)\) &
Derived displacement amplitude &
\tblcell{Raw scalar displacement amplitude obtained from the modulus of the displacement projection:
\[
A(x)=|\Phi_{\rm disp}(x)|.
\]} \\

\(Z_{\rm disp}\) &
Derived normalization &
\tblcell{Spectral normalization of the internal displacement profile:
\[
Z_{\rm disp}=\langle\Xi_{\rm disp}|\Xi_{\rm disp}\rangle_P.
\]} \\

\(A_{\rm phys}(x)\) &
Derived physical field &
\tblcell{Canonically normalized observable displacement amplitude:
\[
A_{\rm phys}(x)=\sqrt{Z_{\rm disp}}\,A(x).
\]} \\

\(A_0^2\) &
Derived vacuum quantity &
\tblcell{Stable nonzero displacement vacuum from the bounded displacement potential \(V_{\rm disp}(A)=V_0+\alpha_AA^2+\beta_AA^4\):
\[
A_0^2=-\frac{\alpha_A}{2\beta_A}.
\]} \\

\(v_A\) &
Derived vacuum quantity &
\tblcell{Canonically normalized displacement vacuum:
\[
v_A=\sqrt{Z_{\rm disp}}A_0.
\]} \\

\(m_A^2\) &
Derived curvature mode &
\tblcell{Positive displacement-curvature scale about the stable nonzero displacement vacuum:
\[
m_A^2=
\left.\frac{d^2V_{\rm disp}}{dA^2}\right|_{A_0}
=-4\alpha_A=8\beta_AA_0^2.
\]} \\

\(\mathcal M_{\rm eff}\) &
Derived inertia scale &
\tblcell{Effective matter inertia generated by displacement vacuum weighted by internal matter--displacement overlap:
\[
\mathcal M_{\rm eff}=g_0\mathcal I_AA_0.
\]} \\

\(T_{\mu\nu}^{({\rm disp})}\) &
Derived source tensor &
\tblcell{Displacement-sector stress-energy contribution sourcing the gravitational stiffness projection.} \\

\(P_A[\Psi]\) &
Derived projection sector &
\tblcell{Compact phase projection generating the observable electromagnetic gauge connection:
\[
P_A[\Psi]\rightarrow\theta(x)\rightarrow\mathcal A_\mu(x).
\]} \\

\(X_\mu\) &
Derived gauge-invariant structure &
\tblcell{Compact phase-gradient combination invariant under local phase reparameterization:
\[
X_\mu=\partial_\mu\theta-q\mathcal A_\mu.
\]} \\

\(F_{\mu\nu}\) &
Derived low-energy field strength &
\tblcell{Observable electromagnetic field strength emerging in the low-energy compact phase limit:
\[
F_{\mu\nu}=\partial_\mu\mathcal A_\nu-\partial_\nu\mathcal A_\mu.
\]} \\

\(T_A,\,T_X\) &
Derived trace split &
\tblcell{Compact/spatial trace fractions of the \(3+1\) projection:
\[
T_A=\frac14,
\qquad
T_X=\frac34.
\]
These are fixed by the projection trace, not by calibration.} \\

\(\Pi_{13}\) &
Derived finite-rank slice &
\tblcell{Finite odd spatial-trace slice selected by the fundamental compact electromagnetic winding:
\[
\Pi_{13}=|\phi_1\rangle\langle\phi_1|+|\phi_3\rangle\langle\phi_3|.
\]
The \(n=3\) channel is the dominant spatial-trace radial lock; the \(n=1\) channel is the parity-compatible coherence tail.} \\

\(K_X\) &
Derived normalized stiffness operator &
\tblcell{Dimensionless radial stiffness operator normalized by the first radial spectral gap:
\[
K_X=\frac{O_X-\lambda_0}{\lambda_1-\lambda_0},
\qquad
\mu_{\min}^2=\lambda_1-\lambda_0.
\]} \\

\(M_A\) &
Derived compact-return matrix &
\tblcell{Finite compact-return transfer matrix generated by the primitive compact trace returns \(A^3\) and \(A^4\):
\[
M_A=
\begin{pmatrix}
T_A^3&T_A^4\\
1&0
\end{pmatrix}.
\]} \\

\(D_A\) &
Derived Fredholm determinant &
\tblcell{Compact-return Fredholm determinant resumming recurrent interior compact paths:
\[
D_A=\det(I-M_A)=1-T_A^3-T_A^4.
\]
For \(T_A=1/4\), \(D_A=251/256\).} \\

\(\operatorname{Adm}_{\Pi_{13}}\) &
Derived admissible boundary words &
\tblcell{Finite-rank terminal word set of the compact boundary map:
\[
\operatorname{Adm}_{\Pi_{13}}
=
\{\epsilon,\ A^3,\ A^4,\ XA^5,\ -XA^8\}.
\]
The terms represent direct passage, primitive compact returns, spatially weighted boundary leakage, and double full-trace exclusion required by finite-rank projection consistency.} \\

\(N_{\rm bc}\) &
Derived boundary cofactor &
\tblcell{Boundary-orientation cofactor obtained from the admissible terminal word set:
\[
N_{\rm bc}=1+T_A^3+T_A^4+T_X(T_A^5-T_A^8).
\]
This is not a fitted correction; it is the signed boundary cofactor of the finite-rank compact return map.} \\

\(c_{\rm bc}\) &
Derived compact boundary constant &
\tblcell{Compact boundary constant produced by the finite-rank Fredholm determinant and boundary cofactor:
\[
\begin{aligned}
c_{\rm bc}
&=
T_X\left[
1+T_A^2
-
T_A^6\frac{N_{\rm bc}}{D_A}
\right],\\
c_{\rm bc}
&=
\frac{3354902985}{4211081216}
=
0.7966844648478991.
\end{aligned}
\]} \\

\(P_{\rm geom}^{\rm bc}\) &
Derived finite-rank projection operator &
\tblcell{Compact boundary-orientation projection operator acting on the finite odd spatial slice:
\[
P_{\rm geom}^{\rm bc}
=
\Pi_{13}K_X^2
\left(
c_{\rm bc}+w^2+\frac34w^4
\right)
K_X^2\Pi_{13}.
\]
This replaces scalar-kernel treatments of the compact projection and preserves the finite-rank orthogonal-slice structure.} \\

\(p_1^{\rm bc}\) &
Numerically evaluated derived output &
\tblcell{Parity-compatible coherence-tail probability produced by \(P_{\rm geom}^{\rm bc}\):
\[
p_1^{\rm bc}=7.6528903366\times10^{-4}.
\]
This lowers the dominant \(n=3\) radial lock to the compact boundary-resolved value.} \\

\(R_{A,\star}^{-2}=Z_A\) &
Derived compact stiffness &
\tblcell{Stationary compact phase stiffness generated by the finite-rank compact boundary map. This is the compact electromagnetic normalization entering \(\alpha_{\rm PR}^{-1}=4\pi Z_A\).} \\

\(C_A\) &
Derived compact/radial closure ratio &
\tblcell{Macroscopic compact-to-radial closure ratio fixed by compact/radial stationarity:
\[
C_A=\frac{R_{A,\star}^{-2}}{\mu_{\min}^2}.
\]
It is not independent of the compact coherence expectation.} \\

\(\kappa_{\rm eff}\) &
Derived under boundary-orientation closure &
\tblcell{Microscopic compact coherence expectation explaining the macroscopic closure ratio:
\[
\kappa_{\rm eff}
=
\langle\Omega_A|\widehat{\kappa}_A|\Omega_A\rangle_P
=
C_A.
\]
Under the finite-rank boundary-orientation lemma, stationarity locks \(\kappa_{\rm eff}=C_A\).} \\

\(\alpha_{{\rm PR},bc}^{-1}\) &
Derived compact-boundary output &
\tblcell{Fine-structure normalization produced by the compact boundary map. The low-energy laboratory value is not used as an input:
\[
\alpha_{{\rm PR},bc}^{-1}
=
4\pi R_{A,\star}^{-2}
=
137.036361812007,
\]
at current nonperturbative precision.} \\

\(P_X[\Psi]\) &
Derived response sector &
\tblcell{Projection response sector encoding internal spectral self-energy corrections:
\[
P_X[\Psi]\rightarrow\Sigma_X.
\]} \\

\(\rho_X(\mu^2)\) &
Derived spectral density &
\tblcell{Internal excitation density with support only above the radial gap \(\mu_{\min}^2\).} \\

\(F(k^2)\) &
Derived response kernel &
\tblcell{Subtracted projection self-energy kernel satisfying \(F(0)=0\), preserving the massless low-energy spin--2 pole.} \\

\(F_E(k^2)\) &
Derived residual kernel &
\tblcell{Newton-normalized residual projection kernel obeying
\[
F_E(k^2)=O(k^4)
\]
in the low-energy exterior limit.} \\

\(\widehat{\Sigma}_{X}^{\rm GW}\) &
Derived GW response &
\tblcell{Projection self-energy correction acting on the Kerr/Teukolsky gravitational-wave operator:
\[
\widehat{\Sigma}_{X}^{\rm GW}
=
-F_E(O_{\rm Teuk}^{(s)}).
\]} \\

\(P_\theta[\Psi]\) &
Derived homogeneous response &
\tblcell{Homogeneous compact phase response sector used for large-scale projection-energy bookkeeping:
\[
P_\theta[\Psi]\rightarrow H_{\rm eff}(t).
\]} \\

\end{longtable}

\endgroup

\section{Gravitational Projection Chain}
\label{sec:gravitational_projection_chain}

The gravitational sector is the stiffness projection of the master field. This
appendix collects the chain by which the internal spectral amplitudes generate
the metric seed tensor, the determinant-normalized effective metric, the
trace-free weak-field Einstein equation, the finite-core saturation scale, and
the Kerr exterior recovery. The purpose is to provide the derivation ledger for
the gravitational equations used in the main text, not to introduce an
independent gravitational postulate.

\subsection{Internal Mode Amplitudes}

The master projection field is defined over observable spacetime and the
internal spectral manifold,
\begin{equation}
\Psi=\Psi(x,\xi).
\end{equation}
Using the internal spectral basis, the field expands as
\begin{equation}
\boxed{
\Psi(x,\xi)
=
\sum_{n,m}
c_{n,m}(x)U_{n,m}(\xi).
}
\end{equation}
The internal modes satisfy
\begin{equation}
\boxed{
O_{\rm int}U_{n,m}
=
\Lambda_{n,m}U_{n,m},
}
\end{equation}
with orthonormality
\begin{equation}
\boxed{
\int_{\mathcal M_{\rm int}}
d\mu_{\rm int}\,
U_{j,k}^{*}(\xi)U_{n,m}(\xi)
=
\delta_{jn}\delta_{km}.
}
\end{equation}
The observable spacetime dependence is carried by the coefficients
\(c_{n,m}(x)\), while the internal stiffness data are carried by
\(\Lambda_{n,m}\).

\subsection{Metric Seed Tensor}

The gravitational projection begins with the symmetric bilinear seed tensor
\begin{equation}
\boxed{
Q_{\mu\nu}(x)
=
\int_{\mathcal M_{\rm int}}
d\mu_{\rm int}
\left[
\partial_\mu\Psi^\dagger
\partial_\nu\Psi
+
\partial_\nu\Psi^\dagger
\partial_\mu\Psi
\right].
}
\end{equation}
This tensor is manifestly symmetric:
\begin{equation}
\boxed{
Q_{\mu\nu}=Q_{\nu\mu}.
}
\end{equation}

Substituting the spectral expansion gives
\begin{equation}
\partial_\mu\Psi
=
\sum_{n,m}
(\partial_\mu c_{n,m})U_{n,m},
\qquad
\partial_\mu U_{n,m}(\xi)=0.
\end{equation}
The conjugate derivative is
\begin{equation}
\partial_\mu\Psi^\dagger
=
\sum_{j,k}
(\partial_\mu c_{j,k}^{*})U_{j,k}^{*}.
\end{equation}
Using the orthonormality relation, the internal integral collapses:
\begin{equation}
\int d\mu_{\rm int}\,
\partial_\mu\Psi^\dagger\partial_\nu\Psi
=
\sum_{n,m}
(\partial_\mu c_{n,m}^{*})
(\partial_\nu c_{n,m}).
\end{equation}
The seed tensor is therefore
\begin{equation}
\boxed{
Q_{\mu\nu}
=
\sum_{n,m}
\left[
\partial_\mu c_{n,m}^{*}
\partial_\nu c_{n,m}
+
\partial_\nu c_{n,m}^{*}
\partial_\mu c_{n,m}
\right]
=
2\,{\rm Re}
\left[
\sum_{n,m}
\partial_\mu c_{n,m}^{*}
\partial_\nu c_{n,m}
\right].
}
\end{equation}
This is the observable stiffness tensor of the projected spectral amplitudes.

\subsection{Determinant-Normalized Effective Metric}

The seed tensor still contains an arbitrary conformal scale. The projected
metric removes that scale by determinant normalization:
\begin{equation}
\boxed{
g_{\mu\nu}^{\rm eff}
=
\frac{Q_{\mu\nu}}{|\det Q|^{1/4}}.
}
\end{equation}
The exponent \(1/4\) is fixed by the four observable spacetime dimensions. If
\(Q_{\mu\nu}\rightarrow sQ_{\mu\nu}\), then:
\begin{equation}
\det(sQ)=s^4\det Q,
\qquad
|\det(sQ)|^{1/4}=s|\det Q|^{1/4},
\end{equation}
so the normalized metric is invariant:
\begin{equation}
\boxed{
Q_{\mu\nu}\rightarrow sQ_{\mu\nu}
\quad
\Longrightarrow
\quad
g_{\mu\nu}^{\rm eff}\rightarrow g_{\mu\nu}^{\rm eff}.
}
\end{equation}
Its determinant is
\begin{equation}
\det(g_{\mu\nu}^{\rm eff})
=
\frac{\det Q}{|\det Q|}
=
{\rm sgn}(\det Q).
\end{equation}
For Lorentzian signature,
\begin{equation}
\boxed{
\det(g_{\mu\nu}^{\rm eff})=-1.
}
\end{equation}
The determinant-normalized projection therefore removes the local volume-scale
degree of freedom while preserving the Lorentzian causal structure.

\subsection{Trace-Free Weak-Field Limit}

In the weak-field regime, write
\begin{equation}
Q_{\mu\nu}
=
\eta_{\mu\nu}+q_{\mu\nu},
\qquad
|q_{\mu\nu}|\ll1,
\end{equation}
and define
\begin{equation}
q=\eta^{\mu\nu}q_{\mu\nu}.
\end{equation}
To first order,
\begin{equation}
|\det Q|^{-1/4}
=
1-\frac14 q+O(q^2).
\end{equation}
The determinant-normalized metric becomes
\begin{equation}
g_{\mu\nu}^{\rm eff}
=
\eta_{\mu\nu}
+
q_{\mu\nu}
-
\frac14\eta_{\mu\nu}q
+
O(q^2).
\end{equation}
Thus the physical perturbation is
\begin{equation}
\boxed{
h_{\mu\nu}
=
q_{\mu\nu}
-
\frac14\eta_{\mu\nu}q.
}
\end{equation}
Taking the trace gives
\begin{equation}
\boxed{
h=\eta^{\mu\nu}h_{\mu\nu}=0.
}
\end{equation}

In Lorenz gauge,
\begin{equation}
\partial^\mu h_{\mu\nu}=0,
\end{equation}
the linearized Ricci tensor reduces to
\begin{equation}
R_{\mu\nu}^{(1)}
=
-\frac12\Box h_{\mu\nu}.
\end{equation}
Because the determinant-normalized perturbation is trace-free, the directly
projected weak-field equation couples to the trace-free source combination:
\begin{equation}
\boxed{
R_{\mu\nu}^{(1)}
-
\frac14\eta_{\mu\nu}R^{(1)}
=
\frac{8\pi G_N}{c^4}
\left(
T_{\mu\nu}
-
\frac14\eta_{\mu\nu}T
\right).
}
\end{equation}
In trace-free Lorenz gauge this becomes
\begin{equation}
\boxed{
-\frac12\Box h_{\mu\nu}
=
\frac{8\pi G_N}{c^4}
\left(
T_{\mu\nu}
-
\frac14\eta_{\mu\nu}T
\right).
}
\end{equation}
Both sides are trace-free. Stress-energy conservation and the contracted Bianchi
identity restore the standard Einstein equation up to the usual cosmological
integration constant:
\begin{equation}
\boxed{
G_{\mu\nu}
+
\Lambda g_{\mu\nu}
=
\frac{8\pi G_N}{c^4}T_{\mu\nu}.
}
\end{equation}
The determinant-normalized gravitational projection therefore recovers the
Einstein limit through the unimodular trace-free form.

\subsection{Finite-Core Saturation}

The radial projection branch has a nonzero first spectral gap,
\begin{equation}
\boxed{
\mu_{\min}^2
=
\lambda_1-\lambda_0
=
3.052966743096.
}
\end{equation}
This gap sets a finite curvature ceiling,
\begin{equation}
\boxed{
R_{\max}
=
\chi_R\Lambda_X^2\mu_{\min}^2
<
\infty,
}
\end{equation}
where \(\chi_R\) is the fixed projection conversion from internal stiffness units
to macroscopic curvature units.

Matching the exterior Schwarzschild curvature scale to the projection ceiling,
\begin{equation}
R_{\rm GR}(r)
=
\frac{G_NM}{c^2r^3}
=
R_{\max},
\end{equation}
gives the finite-core radius
\begin{equation}
\boxed{
r_c(M)
=
\left(
\frac{G_NM}{c^2R_{\max}}
\right)^{1/3}.
}
\end{equation}
Since \(R_{\max}<\infty\), the core radius is strictly positive:
\begin{equation}
\boxed{
r_c>0.
}
\end{equation}

A minimal regularized mass profile satisfying exterior recovery and central
volume scaling is
\begin{equation}
\boxed{
m_{\rm PR}(r)
=
M
\frac{r^3}{r^3+r_c^3}.
}
\end{equation}
The central density is finite:
\begin{equation}
\boxed{
\rho_{\rm eff}(0)
=
\frac{3M}{4\pi r_c^3}
=
\frac{3c^2R_{\max}}{4\pi G_N}
\equiv
\rho_{\max}
<
\infty.
}
\end{equation}
The corresponding core curvature invariant is bounded:
\begin{equation}
\boxed{
K_{\rm core}
=
96R_{\max}^2
<
\infty.
}
\end{equation}
The finite-core regularization chain is
\begin{equation}
\boxed{
\mu_{\min}^2>0
\quad\Longrightarrow\quad
R_{\max}<\infty
\quad\Longrightarrow\quad
r_c>0
\quad\Longrightarrow\quad
\rho_{\max}<\infty
\quad\Longrightarrow\quad
K_{\rm PR}(0)<\infty.
}
\end{equation}

\subsection{Exterior Kerr Recovery}

Outside the finite projection core,
\begin{equation}
r>r_c,
\end{equation}
the ordinary matter source vanishes in the exterior vacuum region:
\begin{equation}
T_{\mu\nu}=0.
\end{equation}
The projection correction is also suppressed in the low-energy exterior:
\begin{equation}
\Delta_{\mu\nu}^{(P)}\rightarrow0.
\end{equation}
The exterior field equation reduces to
\begin{equation}
G_{\mu\nu}[g^{\rm eff}]=0.
\end{equation}
Taking the trace gives
\begin{equation}
R=0,
\end{equation}
and substituting back gives
\begin{equation}
\boxed{
R_{\mu\nu}=0,
\qquad
r>r_c.
}
\end{equation}

For stationary, axisymmetric, asymptotically flat exterior boundary conditions
with mass \(M\) and angular momentum \(J\), the vacuum solution is Kerr. Define
\begin{equation}
\mathcal M
=
\frac{G_NM}{c^2},
\qquad
a
=
\frac{J}{Mc}.
\end{equation}
The Kerr structure functions are
\begin{equation}
\Sigma_K
=
r^2+a^2\cos^2\theta,
\qquad
\Delta_K
=
r^2-2\mathcal M r+a^2.
\end{equation}
Projection Relativity allows finite-core corrections only in the interior:
\begin{equation}
\Sigma_{\rm PR}
=
\Sigma_K+\delta\Sigma,
\qquad
\Delta_{\rm PR}
=
\Delta_K+\delta\Delta.
\end{equation}
Exterior recovery imposes
\begin{equation}
\boxed{
\delta\Sigma=0,
\qquad
\delta\Delta=0,
\qquad
r>r_c.
}
\end{equation}
Therefore,
\begin{equation}
\boxed{
g_{\mu\nu}^{\rm PR}
=
g_{\mu\nu}^{\rm Kerr},
\qquad
r>r_c.
}
\end{equation}

\subsection{Closed Gravitational Projection Sequence}

The gravitational projection chain is
\begin{equation}
\boxed{
\Psi
\rightarrow
\{c_{n,m},U_{n,m}\}
\rightarrow
Q_{\mu\nu}
\rightarrow
g_{\mu\nu}^{\rm eff}
\rightarrow
\left[
R_{\mu\nu}^{(1)}
-
\frac14\eta_{\mu\nu}R^{(1)}
=
\frac{8\pi G_N}{c^4}
\left(
T_{\mu\nu}
-
\frac14\eta_{\mu\nu}T
\right)
\right].
}
\end{equation}
The Bianchi-completed macroscopic Einstein limit is
\begin{equation}
\boxed{
G_{\mu\nu}
+
\Lambda g_{\mu\nu}
=
\frac{8\pi G_N}{c^4}T_{\mu\nu}.
}
\end{equation}
The finite-core chain is
\begin{equation}
\boxed{
\mu_{\min}^2
\rightarrow
R_{\max}
\rightarrow
r_c(M)
=
\left(
\frac{G_NM}{c^2R_{\max}}
\right)^{1/3}
\rightarrow
K_{\rm PR}(0)<\infty.
}
\end{equation}
The exterior recovery chain is
\begin{equation}
\boxed{
r>r_c
\rightarrow
T_{\mu\nu}=0
\rightarrow
\Delta_{\mu\nu}^{(P)}\rightarrow0
\rightarrow
R_{\mu\nu}=0
\rightarrow
g_{\mu\nu}^{\rm PR}
=
g_{\mu\nu}^{\rm Kerr}.
}
\end{equation}

This completes the gravitational projection chain. The same master-field
stiffness projection recovers the weak-field Einstein limit, replaces the
classical singular endpoint with a finite core, and preserves the Kerr exterior
outside the projection-saturated region.

\section{Displacement Projection Chain}
\label{sec:displacement_projection_chain}

The displacement sector is the scalar-amplitude projection of the master field.
This appendix collects the chain by which the master projection state generates
an observable displacement amplitude, a stable nonzero displacement vacuum, a
positive displacement-curvature mode, and an effective matter inertia scale. The
purpose is to record the derivation ledger for the displacement equations used
in the main text, not to introduce an independent scalar sector.

\subsection{Master-Field Projection}

At a fixed observable spacetime point \(x\), the master field is written in the
internal spectral basis as
\begin{equation}
\boxed{
|\Psi_x\rangle
=
\sum_{n,m}
c_{n,m}(x)|U_{n,m}\rangle .
}
\end{equation}
Equivalently, in coordinate form,
\begin{equation}
\boxed{
\Psi(x,\xi)
=
\sum_{n,m}
c_{n,m}(x)U_{n,m}(\xi).
}
\end{equation}
The internal basis is orthonormal:
\begin{equation}
\boxed{
\langle U_{j,k}|U_{n,m}\rangle_P
=
\delta_{jn}\delta_{km}.
}
\end{equation}

Let \(U_{\rm disp}(\xi)\) denote the normalized internal displacement direction,
\begin{equation}
\boxed{
\langle U_{\rm disp}|U_{\rm disp}\rangle_P=1.
}
\end{equation}
The complex displacement projection is
\begin{equation}
\boxed{
\Phi_{\rm disp}(x)
=
\langle U_{\rm disp}|\Psi_x\rangle_P.
}
\end{equation}
In integral form,
\begin{equation}
\boxed{
\Phi_{\rm disp}(x)
=
\int_{\mathcal M_{\rm int}}
d\mu_{\rm int}\,
U_{\rm disp}^{*}(\xi)\Psi(x,\xi).
}
\end{equation}
Substituting the spectral expansion gives
\begin{equation}
\boxed{
\Phi_{\rm disp}(x)
=
\sum_{n,m}
d_{n,m}^{({\rm disp})}c_{n,m}(x),
\qquad
d_{n,m}^{({\rm disp})}
=
\langle U_{\rm disp}|U_{n,m}\rangle_P.
}
\end{equation}
The observable displacement amplitude is the modulus of this projected field:
\begin{equation}
\boxed{
A(x)
=
|\Phi_{\rm disp}(x)|,
\qquad
A^2(x)=\Phi_{\rm disp}^{\dagger}(x)\Phi_{\rm disp}(x).
}
\end{equation}
Thus the displacement projection is
\begin{equation}
\boxed{
P_{\rm disp}[\Psi](x)=A(x).
}
\end{equation}

\subsection{Rank-One Projection Form}

The same projection can be written as a rank-one internal projector,
\begin{equation}
\boxed{
\widehat{\Pi}_{\rm disp}
=
|U_{\rm disp}\rangle\langle U_{\rm disp}|.
}
\end{equation}
Acting on the master state,
\begin{equation}
\widehat{\Pi}_{\rm disp}|\Psi_x\rangle
=
|U_{\rm disp}\rangle
\langle U_{\rm disp}|\Psi_x\rangle_P
=
\Phi_{\rm disp}(x)|U_{\rm disp}\rangle .
\end{equation}
Its internal norm is
\begin{equation}
\left\|
\widehat{\Pi}_{\rm disp}|\Psi_x\rangle
\right\|_P^2
=
\Phi_{\rm disp}^{\dagger}(x)\Phi_{\rm disp}(x)
\langle U_{\rm disp}|U_{\rm disp}\rangle_P
=
A^2(x).
\end{equation}
Therefore,
\begin{equation}
\boxed{
A(x)
=
\left\|
\widehat{\Pi}_{\rm disp}|\Psi_x\rangle
\right\|_P.
}
\end{equation}
This form makes explicit that the displacement amplitude is a projection norm,
not an independent inserted scalar field.

\subsection{Amplitude--Phase Decomposition}

For nonzero projected displacement field, write
\begin{equation}
\boxed{
\Phi_{\rm disp}(x)
=
A(x)e^{i\theta(x)}.
}
\end{equation}
The compact phase obeys the same winding identification used in the compact
electromagnetic sector:
\begin{equation}
\boxed{
\theta\sim\theta+2\pi n,
\qquad
n\in\mathbb Z.
}
\end{equation}

The ordinary derivative gives
\begin{equation}
\partial_\mu\Phi_{\rm disp}
=
e^{i\theta}
\left(
\partial_\mu A
+
iA\partial_\mu\theta
\right),
\end{equation}
and
\begin{equation}
\partial_\mu\Phi_{\rm disp}^{\dagger}
=
e^{-i\theta}
\left(
\partial_\mu A
-
iA\partial_\mu\theta
\right).
\end{equation}
The kinetic bilinear separates as
\begin{align}
\partial_\mu\Phi_{\rm disp}^{\dagger}
\partial^\mu\Phi_{\rm disp}
&=
(\partial_\mu A)(\partial^\mu A)
+
iA(\partial_\mu A)(\partial^\mu\theta)
-
iA(\partial_\mu\theta)(\partial^\mu A)
+
A^2(\partial_\mu\theta)(\partial^\mu\theta)
\\
&=
(\partial_\mu A)(\partial^\mu A)
+
A^2(\partial_\mu\theta)(\partial^\mu\theta).
\end{align}
Thus,
\begin{equation}
\boxed{
\partial_\mu\Phi_{\rm disp}^{\dagger}
\partial^\mu\Phi_{\rm disp}
=
(\partial_\mu A)(\partial^\mu A)
+
A^2(\partial_\mu\theta)(\partial^\mu\theta).
}
\end{equation}
The first term is the displacement-amplitude kinetic term. The second is the
compact phase kinetic term weighted by the displacement amplitude.

With the compact gauge connection included, define the gauge-invariant phase
gradient
\begin{equation}
\boxed{
X_\mu
=
\partial_\mu\theta-q\mathcal A_\mu.
}
\end{equation}
The covariant amplitude--phase split becomes
\begin{equation}
\boxed{
g_{\rm eff}^{\mu\nu}
\left(D_\mu\Phi_{\rm disp}\right)^\dagger
D_\nu\Phi_{\rm disp}
=
(\partial_\mu A)(\partial^\mu A)
+
A^2X_\mu X^\mu .
}
\end{equation}

\subsection{Displacement-Sector Action}

After projecting the master action onto the displacement amplitude, the
displacement-sector action is
\begin{equation}
\boxed{
S_{\rm disp}
=
\int d^4x\sqrt{-g_{\rm eff}}\,
Z_{\rm disp}
\left[
-\frac12
(\partial_\mu A)(\partial^\mu A)
-
V_{\rm disp}(A)
\right].
}
\end{equation}
The minimal bounded displacement potential is
\begin{equation}
\boxed{
V_{\rm disp}(A)
=
V_0+\alpha_AA^2+\beta_AA^4,
\qquad
\beta_A>0.
}
\end{equation}
The displacement equation of motion is
\begin{equation}
\boxed{
\Box A
-
2\alpha_AA
-
4\beta_AA^3
=
0.
}
\end{equation}

\subsection{Stable Nonzero Displacement Vacuum}

Stationarity of the displacement potential gives
\begin{equation}
\frac{dV_{\rm disp}}{dA}
=
2\alpha_AA+4\beta_AA^3
=
2A
\left(
\alpha_A+2\beta_AA^2
\right).
\end{equation}
The nonzero branch satisfies
\begin{equation}
\boxed{
A_0^2
=
-\frac{\alpha_A}{2\beta_A}.
}
\end{equation}
A real stable nonzero displacement vacuum requires
\begin{equation}
\boxed{
\alpha_A<0,
\qquad
\beta_A>0.
}
\end{equation}
Then
\begin{equation}
\boxed{
A_0
=
\sqrt{
-\frac{\alpha_A}{2\beta_A}
}.
}
\end{equation}
The canonically normalized displacement vacuum is
\begin{equation}
\boxed{
v_A
=
\sqrt{Z_{\rm disp}}\,A_0.
}
\end{equation}

\subsection{Displacement Fluctuation and Curvature Scale}

Let
\begin{equation}
\boxed{
A(x)=A_0+\eta_A(x).
}
\end{equation}
The displacement-curvature scale is
\begin{equation}
m_A^2
=
\left.
\frac{d^2V_{\rm disp}}{dA^2}
\right|_{A_0}.
\end{equation}
Since
\begin{equation}
\frac{d^2V_{\rm disp}}{dA^2}
=
2\alpha_A+12\beta_AA^2,
\end{equation}
substitution of \(A_0^2=-\alpha_A/(2\beta_A)\) gives
\begin{equation}
\boxed{
m_A^2
=
-4\alpha_A
=
8\beta_AA_0^2.
}
\end{equation}
Because \(\alpha_A<0\),
\begin{equation}
\boxed{
m_A^2>0.
}
\end{equation}
The quadratic displacement fluctuation Lagrangian is
\begin{equation}
\boxed{
\mathcal L_{\rm disp}^{(2)}
=
-\frac12
(\partial_\mu\eta_A)(\partial^\mu\eta_A)
-
\frac12m_A^2\eta_A^2.
}
\end{equation}
The corresponding fluctuation equation is
\begin{equation}
\boxed{
(\Box-m_A^2)\eta_A=0.
}
\end{equation}

\subsection{Spectral Normalization}

The internal displacement profile has spectral normalization
\begin{equation}
\boxed{
Z_{\rm disp}
=
\int_{\mathcal M_{\rm int}}
d\mu_{\rm int}
|\Xi_{\rm disp}(\xi)|^2
=
\sum_{n,m}
|\zeta_{n,m}^{({\rm disp})}|^2.
}
\end{equation}
The canonically normalized displacement amplitude is
\begin{equation}
\boxed{
A_{\rm phys}(x)
=
\sqrt{Z_{\rm disp}}\,A(x).
}
\end{equation}
The canonically normalized displacement vacuum is
\begin{equation}
\boxed{
v_A
=
\sqrt{Z_{\rm disp}}\,A_0.
}
\end{equation}
The canonical fluctuation is
\begin{equation}
\boxed{
a_{\rm phys}(x)
=
\sqrt{Z_{\rm disp}}\,\eta_A(x).
}
\end{equation}
Therefore,
\begin{equation}
\boxed{
A_{\rm phys}(x)
=
v_A+a_{\rm phys}(x).
}
\end{equation}

\subsection{Effective Matter Inertia}

Matter inertia is generated through the matter--displacement overlap interaction,
written schematically as
\begin{equation}
\boxed{
\mathcal L_{\rm ov}
=
-g_{\rm eff}\mathcal O_MA_{\rm phys}.
}
\end{equation}
Using
\begin{equation}
A_{\rm phys}=v_A+a_{\rm phys},
\end{equation}
the overlap interaction becomes
\begin{equation}
\mathcal L_{\rm ov}
=
-g_{\rm eff}v_A\mathcal O_M
-
g_{\rm eff}a_{\rm phys}\mathcal O_M.
\end{equation}
The first term is the effective matter-inertia term. Hence
\begin{equation}
\boxed{
\mathcal M_{\rm eff}
=
g_{\rm eff}v_A.
}
\end{equation}

The effective overlap coefficient is determined by the internal displacement
overlap,
\begin{equation}
\boxed{
g_{\rm eff}
=
\frac{g_0\mathcal I_A}{\sqrt{Z_{\rm disp}}}.
}
\end{equation}
Therefore,
\begin{equation}
\boxed{
\mathcal M_{\rm eff}
=
g_{\rm eff}v_A
=
\frac{g_0\mathcal I_A}{\sqrt{Z_{\rm disp}}}
\left(
\sqrt{Z_{\rm disp}}A_0
\right)
=
g_0\mathcal I_AA_0.
}
\end{equation}
The physical effective inertia is independent of an arbitrary normalization of
the internal displacement profile. For multiple matter projection profiles, the overlap generalizes to the matrix
form
\begin{equation}
\boxed{
g_{ab}^{\rm eff}
=
\frac{g_0I_{ab}^{(A)}}{\sqrt{Z_{\rm disp}}},
\qquad
\mathcal M_{ab}^{\rm eff}
=
g_{ab}^{\rm eff}v_A
=
g_0I_{ab}^{(A)}A_0.
}
\end{equation}

\subsection{Displacement Stress-Energy and Gravitational Source Response}

The displacement stress-energy tensor follows from the displacement-sector
action:
\begin{equation}
\boxed{
T_{\mu\nu}^{({\rm disp})}
=
Z_{\rm disp}
\left[
\partial_\mu A\,\partial_\nu A
-
g_{\mu\nu}^{\rm eff}
\left(
\frac12
(\partial_\alpha A)(\partial^\alpha A)
+
V_{\rm disp}(A)
\right)
\right].
}
\end{equation}
The displacement vacuum contributes to the physical matter source through the
effective inertia scale \(\mathcal M_{\rm eff}\). In the determinant-normalized
weak-field limit, this source enters the gravitational stiffness projection
through the trace-free unimodular combination
\begin{equation}
\boxed{
T_{\mu\nu}^{({\rm disp})}
-
\frac14
g_{\mu\nu}^{\rm eff}
T^{({\rm disp})}.
}
\end{equation}
Stress-energy conservation and the contracted Bianchi identity restore the scalar
trace through the usual cosmological integration constant, yielding the standard
Einstein source equation in the macroscopic low-energy limit. The displacement
sector therefore supplies gravitational source response consistently with the
trace-free weak-field recovery derived in the gravitational projection chain.

\subsection{Displacement--Phase Consistency}

The covariant compact phase contribution generated by the displacement field is
\begin{equation}
\boxed{
A^2X_\mu X^\mu.
}
\end{equation}
Expanding about the displacement vacuum,
\begin{equation}
A=A_0+\eta_A,
\end{equation}
gives
\begin{equation}
\boxed{
A^2X_\mu X^\mu
=
A_0^2X_\mu X^\mu
+
2A_0\eta_A X_\mu X^\mu
+
\eta_A^2X_\mu X^\mu.
}
\end{equation}
The compact phase current is
\begin{equation}
\boxed{
J_\theta^\mu
=
A^2X^\mu.
}
\end{equation}
In the coherent compact-phase vacuum,
\begin{equation}
\boxed{
X_\mu^{\rm vac}=0,
}
\end{equation}
so
\begin{equation}
\boxed{
J_\theta^{\mu,{\rm vac}}=0.
}
\end{equation}
The compact electromagnetic gauge normalization and boundary-resolved closure
are derived in the compact phase projection chain.

\subsection{Stability Closure}

The displacement sector is stable when
\begin{equation}
\boxed{
\beta_A>0,
\qquad
\alpha_A<0,
\qquad
Z_{\rm disp}>0,
\qquad
\mu_{\min}^2>0,
\qquad
X_\mu^{\rm vac}=0.
}
\end{equation}
These conditions imply
\begin{equation}
\boxed{
V_{\rm disp}(A)\ \text{is bounded below},
\qquad
m_A^2>0,
\qquad
Z_{\rm disp}<\infty,
\qquad
J_\theta^{\mu,{\rm vac}}=0.
}
\end{equation}

\subsection{Closed Displacement Projection Sequence}

The displacement projection chain is
\begin{equation}
\boxed{
|\Psi_x\rangle
\rightarrow
\widehat{\Pi}_{\rm disp}|\Psi_x\rangle
\rightarrow
\Phi_{\rm disp}(x)
=
\langle U_{\rm disp}|\Psi_x\rangle_P
\rightarrow
A(x)
=
|\Phi_{\rm disp}(x)|.
}
\end{equation}
The displacement-vacuum chain is
\begin{equation}
\boxed{
A(x)
\rightarrow
V_{\rm disp}(A)
=
V_0+\alpha_AA^2+\beta_AA^4
\rightarrow
A_0^2
=
-\frac{\alpha_A}{2\beta_A}
\rightarrow
A(x)=A_0+\eta_A(x).
}
\end{equation}
The displacement-curvature chain is
\begin{equation}
\boxed{
m_A^2
=
\left.
\frac{d^2V_{\rm disp}}{dA^2}
\right|_{A_0}
=
-4\alpha_A
=
8\beta_AA_0^2.
}
\end{equation}
The spectral-normalization chain is
\begin{equation}
\boxed{
Z_{\rm disp}
=
\int_{\mathcal M_{\rm int}}
d\mu_{\rm int}
|\Xi_{\rm disp}(\xi)|^2,
\qquad
A_{\rm phys}
=
\sqrt{Z_{\rm disp}}\,A,
\qquad
v_A
=
\sqrt{Z_{\rm disp}}\,A_0.
}
\end{equation}
The effective-inertia chain is
\begin{equation}
\boxed{
\mathcal M_{\rm eff}
=
g_{\rm eff}v_A
=
g_0\mathcal I_AA_0.
}
\end{equation}
The displacement--gravity chain is
\begin{equation}
\boxed{
A_0
\rightarrow
\mathcal M_{\rm eff}
\rightarrow
T_{\mu\nu}^{({\rm disp})}
\rightarrow
T_{\mu\nu}^{({\rm disp})}
-
\frac14
g_{\mu\nu}^{\rm eff}
T^{({\rm disp})}
\rightarrow
g_{\mu\nu}^{\rm eff}.
}
\end{equation}
The displacement--phase consistency chain is
\begin{equation}
\boxed{
\Phi_{\rm disp}=Ae^{i\theta}
\rightarrow
g_{\rm eff}^{\mu\nu}
(D_\mu\Phi_{\rm disp})^\dagger D_\nu\Phi_{\rm disp}
=
(\partial_\mu A)(\partial^\mu A)
+
A^2X_\mu X^\mu,
\qquad
X_\mu=\partial_\mu\theta-q\mathcal A_\mu.
}
\end{equation}
The full displacement-sector closure is
\begin{equation}
\boxed{
\Psi
\rightarrow
P_{\rm disp}[\Psi]
\rightarrow
\Phi_{\rm disp}(x)
\rightarrow
A(x)
\rightarrow
A_0
\rightarrow
\eta_A
\rightarrow
m_A
\rightarrow
A_{\rm phys}
\rightarrow
\mathcal M_{\rm eff}
\rightarrow
T_{\mu\nu}^{({\rm disp})}
\rightarrow
g_{\mu\nu}^{\rm eff}.
}
\end{equation}

This completes the displacement projection chain. The displacement sector is not
an independent external scalar field. It is the scalar-amplitude projection of
the master field. Its stable nonzero vacuum generates displacement curvature,
effective matter inertia, gravitational source response, and consistent coupling
to the compact phase sector.

\section{Electromagnetic Projection Chain}
\label{sec:electromagnetic_projection_chain}

The electromagnetic sector is the compact-phase projection of the master field.
It is not introduced as an independent external \(U(1)\) gauge bundle. It is
generated by the compact coordinate \(\theta\) of the same internal
radial--compact manifold used by the gravitational and displacement sectors.
This appendix collects the compact phase chain: the compact operator, winding
basis, gauge-invariant phase gradient, field-strength curvature, kinetic
normalization, compact boundary closure, and magnetic area-law interpretation.

\subsection{Compact Phase Geometry}

The minimal internal manifold is
\begin{equation}
\boxed{
\mathcal M_{\rm int}
=
\mathbb R_w\times S^1_\theta,
\qquad
\theta\sim\theta+2\pi n,
\qquad
n\in\mathbb Z.
}
\end{equation}
The compact coordinate is represented by \(e^{i\theta}\), so the physical
projection is unchanged under
\begin{equation}
\boxed{
e^{i(\theta+2\pi n)}
=
e^{i\theta}.
}
\end{equation}

The internal line element is
\begin{equation}
\boxed{
ds_{\rm int}^2
=
dw^2+R_A^2d\theta^2.
}
\end{equation}
The compact phase measure is
\begin{equation}
\boxed{
d\mu_\theta
=
R_A\,d\theta.
}
\end{equation}

The compact phase Hilbert space is
\begin{equation}
\boxed{
\mathcal H_\theta
=
L^2(S^1_\theta,R_A\,d\theta),
}
\end{equation}
with inner product
\begin{equation}
\boxed{
\langle f|g\rangle_\theta
=
\int_0^{2\pi}
R_A\,d\theta\,
f^*(\theta)g(\theta).
}
\end{equation}

\subsection{Compact Phase Operator and Winding Modes}

The compact phase operator is
\begin{equation}
\boxed{
O_\theta
=
-\frac{1}{R_A^2}
\frac{d^2}{d\theta^2}.
}
\end{equation}
The operator acts on periodic functions,
\begin{equation}
\boxed{
v_m(\theta+2\pi)=v_m(\theta).
}
\end{equation}

Integration by parts on the compact circle gives no boundary contribution, so
\(O_\theta\) is self-adjoint on the periodic domain:
\begin{equation}
\boxed{
\langle f|O_\theta g\rangle_\theta
=
\langle O_\theta f|g\rangle_\theta.
}
\end{equation}

The normalized winding modes are
\begin{equation}
\boxed{
v_m(\theta)
=
\frac{1}{\sqrt{2\pi R_A}}
e^{im\theta},
\qquad
m\in\mathbb Z.
}
\end{equation}
They satisfy
\begin{equation}
\boxed{
\langle v_m|v_n\rangle_\theta
=
\delta_{mn}.
}
\end{equation}

Acting with \(O_\theta\) gives
\begin{equation}
O_\theta v_m
=
-\frac{1}{R_A^2}
\frac{d^2}{d\theta^2}
\left(
\frac{e^{im\theta}}{\sqrt{2\pi R_A}}
\right)
=
\frac{m^2}{R_A^2}
v_m.
\end{equation}
Hence
\begin{equation}
\boxed{
O_\theta v_m
=
\lambda_m^{(\theta)}v_m,
\qquad
\lambda_m^{(\theta)}
=
\frac{m^2}{R_A^2}.
}
\end{equation}

The fundamental electromagnetic winding is \(m=1\). The compact phase stiffness
is therefore
\begin{equation}
\boxed{
Z_A
=
\lambda_1^{(\theta)}
=
R_A^{-2}.
}
\end{equation}

\subsection{Compact Phase Projection}

Let \(U_A(\xi)\) denote the normalized internal compact-phase projection
direction. The compact electromagnetic projection of the master field is
\begin{equation}
\boxed{
\Phi_A(x)
=
\langle U_A|\Psi_x\rangle_P.
}
\end{equation}
In integral form,
\begin{equation}
\boxed{
\Phi_A(x)
=
\int_{\mathcal M_{\rm int}}
d\mu_{\rm int}\,
U_A^*(\xi)\Psi(x,\xi).
}
\end{equation}

The projected compact field is written in amplitude--phase form:
\begin{equation}
\boxed{
\Phi_A(x)
=
A_\theta(x)e^{i\theta(x)}.
}
\end{equation}
Here \(A_\theta(x)=|\Phi_A(x)|\) is the compact projection amplitude, while
\(\theta(x)\) is the observable compact phase.

The compact electromagnetic projection is
\begin{equation}
\boxed{
P_A[\Psi]
\rightarrow
\theta(x)
\rightarrow
\mathcal A_\mu(x).
}
\end{equation}
The gauge connection \(\mathcal A_\mu\) is the observable connection associated
with local compact phase transport.

\subsection{Gauge-Invariant Compact Phase Gradient}

The covariant derivative acting on the compact projected field is
\begin{equation}
\boxed{
D_\mu
=
\partial_\mu-iq\mathcal A_\mu.
}
\end{equation}
Acting on \(\Phi_A=A_\theta e^{i\theta}\),
\begin{align}
D_\mu\Phi_A
&=
(\partial_\mu-iq\mathcal A_\mu)
\left(
A_\theta e^{i\theta}
\right)
\\
&=
e^{i\theta}
\left[
\partial_\mu A_\theta
+
iA_\theta
\left(
\partial_\mu\theta-q\mathcal A_\mu
\right)
\right].
\end{align}

Define the compact phase gradient
\begin{equation}
\boxed{
X_\mu
=
\partial_\mu\theta-q\mathcal A_\mu.
}
\end{equation}
Then
\begin{equation}
\boxed{
D_\mu\Phi_A
=
e^{i\theta}
\left(
\partial_\mu A_\theta+iA_\theta X_\mu
\right).
}
\end{equation}

Under a local compact phase rotation,
\begin{equation}
\boxed{
\theta'
=
\theta+\alpha,
\qquad
\mathcal A_\mu'
=
\mathcal A_\mu
+
q^{-1}\partial_\mu\alpha.
}
\end{equation}
The phase gradient transforms as
\begin{align}
X_\mu'
&=
\partial_\mu\theta'
-
q\mathcal A_\mu'
\\
&=
\partial_\mu(\theta+\alpha)
-
q
\left(
\mathcal A_\mu+q^{-1}\partial_\mu\alpha
\right)
\\
&=
\partial_\mu\theta-q\mathcal A_\mu
=
X_\mu.
\end{align}
Thus
\begin{equation}
\boxed{
X_\mu'=X_\mu.
}
\end{equation}
The compact phase gradient is gauge invariant.

\subsection{Compact Phase Kinetic Split}

The conjugate covariant derivative is
\begin{equation}
\boxed{
(D_\mu\Phi_A)^\dagger
=
e^{-i\theta}
\left(
\partial_\mu A_\theta-iA_\theta X_\mu
\right).
}
\end{equation}
The covariant kinetic bilinear is
\begin{align}
g_{\rm eff}^{\mu\nu}
(D_\mu\Phi_A)^\dagger D_\nu\Phi_A
&=
g_{\rm eff}^{\mu\nu}
\left(
\partial_\mu A_\theta-iA_\theta X_\mu
\right)
\left(
\partial_\nu A_\theta+iA_\theta X_\nu
\right)
\\
&=
(\partial_\mu A_\theta)(\partial^\mu A_\theta)
+
A_\theta^2X_\mu X^\mu.
\end{align}
The imaginary cross terms cancel. Hence
\begin{equation}
\boxed{
g_{\rm eff}^{\mu\nu}
(D_\mu\Phi_A)^\dagger D_\nu\Phi_A
=
(\partial_\mu A_\theta)(\partial^\mu A_\theta)
+
A_\theta^2X_\mu X^\mu.
}
\end{equation}

The first term belongs to the compact projection amplitude. The second term is
the electromagnetic compact phase contribution:
\begin{equation}
\boxed{
P_A[\Psi]
\rightarrow
A_\theta^2X_\mu X^\mu.
}
\end{equation}

The compact phase current is
\begin{equation}
\boxed{
J_\theta^\mu
=
A_\theta^2X^\mu.
}
\end{equation}
The coherent compact phase vacuum satisfies
\begin{equation}
\boxed{
X_\mu^{\rm vac}=0.
}
\end{equation}
Therefore,
\begin{equation}
\boxed{
J_\theta^{\mu,{\rm vac}}=0.
}
\end{equation}
The compact phase vacuum is a zero-current coherent ground state.

\subsection{Gauge Curvature and Field-Strength Invariance}

The electromagnetic field strength is the curvature of the observable compact
phase connection:
\begin{equation}
\boxed{
F_{\mu\nu}
=
\partial_\mu\mathcal A_\nu
-
\partial_\nu\mathcal A_\mu.
}
\end{equation}
Under
\begin{equation}
\mathcal A_\mu'
=
\mathcal A_\mu+q^{-1}\partial_\mu\alpha,
\end{equation}
the curvature becomes
\begin{align}
F_{\mu\nu}'
&=
\partial_\mu\mathcal A_\nu'
-
\partial_\nu\mathcal A_\mu'
\\
&=
\partial_\mu\mathcal A_\nu
-
\partial_\nu\mathcal A_\mu
+
q^{-1}
\left(
\partial_\mu\partial_\nu\alpha
-
\partial_\nu\partial_\mu\alpha
\right)
\\
&=
F_{\mu\nu}.
\end{align}
Thus
\begin{equation}
\boxed{
F_{\mu\nu}'=F_{\mu\nu}.
}
\end{equation}

The Bianchi identity follows from the definition:
\begin{equation}
\boxed{
\nabla_{[\lambda}F_{\mu\nu]}=0.
}
\end{equation}

\subsection{Compact Winding Stiffness and Gauge Kinetic Normalization}

The compact part of the internal kinetic sector is
\begin{equation}
\boxed{
\mathcal L_{\rm int}^{(\theta)}
=
-\frac12
G^{\theta\theta}
D_\theta\Psi^\dagger D_\theta\Psi.
}
\end{equation}
Since
\begin{equation}
\boxed{
G^{\theta\theta}
=
R_A^{-2},
}
\end{equation}
the compact kinetic term is
\begin{equation}
\boxed{
\mathcal L_{\rm int}^{(\theta)}
=
-\frac{1}{2R_A^2}
D_\theta\Psi^\dagger D_\theta\Psi.
}
\end{equation}
For the internal compact eigenproblem,
\begin{equation}
D_\theta\rightarrow\partial_\theta,
\end{equation}
so
\begin{equation}
\boxed{
\mathcal L_{\rm int}^{(\theta)}
=
-\frac{1}{2R_A^2}
\partial_\theta\Psi^\dagger\partial_\theta\Psi.
}
\end{equation}

Projection into observable spacetime maps the compact winding stiffness
\(R_A^{-2}\) into the gauge kinetic normalization. In the normalization used
throughout the paper,
\begin{equation}
\boxed{
\alpha_{\rm PR}^{-1}
=
4\pi Z_A,
\qquad
Z_A=R_A^{-2}.
}
\end{equation}
Equivalently,
\begin{equation}
\boxed{
\alpha_{\rm PR}^{-1}
=
4\pi R_A^{-2}.
}
\end{equation}

This relation fixes the electromagnetic normalization once the compact radius
has been geometrically closed. It does not by itself determine \(R_A\). The
stationary compact radius is fixed by compact/radial coherence and the
finite-rank compact boundary map.

\subsection{Compact/Radial Stationarity}

Let
\begin{equation}
\boxed{
\mu_{\min}^2
=
\lambda_1-\lambda_0
}
\end{equation}
be the radial spectral gap of the projection-trace radial operator. The
macroscopic compact/radial closure ratio is
\begin{equation}
\boxed{
C_A
=
\frac{R_{A,\star}^{-2}}{\mu_{\min}^2}.
}
\end{equation}
The corresponding microscopic compact coherence ratio is
\begin{equation}
\boxed{
\kappa_{\rm eff}
=
\frac{
\langle\Omega_A|\widehat K_X^{(c)}|\Omega_A\rangle_P
}{
\langle\Omega_A|\widehat W_A^{(c)}|\Omega_A\rangle_P
}.
}
\end{equation}
Compact/radial stationarity imposes
\begin{equation}
\boxed{
C_A
=
\kappa_{\rm eff}.
}
\end{equation}
The compact radius is therefore locked by the internal projection geometry:
\begin{equation}
\boxed{
R_{A,\star}^{-2}
=
C_A\mu_{\min}^2.
}
\end{equation}

\subsection{Finite-Rank Compact Boundary Closure}

The compact trace fraction is
\begin{equation}
\boxed{
T_A=\frac14,
}
\end{equation}
and the observable spatial trace fraction is
\begin{equation}
\boxed{
T_X=\frac34.
}
\end{equation}
The interior compact return has primitive words \(A^3\) and \(A^4\), giving
the transfer determinant
\begin{equation}
\boxed{
D_A
=
1-T_A^3-T_A^4.
}
\end{equation}

The finite-rank boundary admits the terminal oriented word set
\begin{equation}
\boxed{
{\rm Adm}_{\Pi_{13}}
=
\{\epsilon,A^3,A^4,XA^5,-XA^8\}.
}
\end{equation}
The boundary cofactor is
\begin{equation}
\boxed{
N_{\rm bc}
=
1+T_A^3+T_A^4+T_X(T_A^5-T_A^8).
}
\end{equation}
The boundary-resolved compact return is
\begin{equation}
\boxed{
R_{\rm bc}
=
\frac{N_{\rm bc}}{D_A}.
}
\end{equation}
The compact boundary constant is
\begin{equation}
\boxed{
c_{\rm bc}
=
T_X
\left[
1+T_A^2-T_A^6R_{\rm bc}
\right].
}
\end{equation}
Equivalently,
\begin{equation}
\boxed{
c_{\rm bc}
=
T_X
\left[
1+T_A^2
-
T_A^6
\frac{
1+T_A^3+T_A^4+T_X(T_A^5-T_A^8)
}{
1-T_A^3-T_A^4
}
\right].
}
\end{equation}
Substitution of
\begin{equation}
T_A=\frac14,
\qquad
T_X=\frac34,
\end{equation}
gives
\begin{equation}
\boxed{
c_{\rm bc}
=
\frac{3354902985}{4211081216}
=
0.7966844648478991.
}
\end{equation}
The compact boundary constant is a terminal-word cofactor of the finite-rank
return map. It is not a fitted correction.

\subsection{Boundary-Resolved Compact Electromagnetic Normalization}

The finite-rank compact projection operator is
\begin{equation}
\boxed{
P_{\rm geom}^{\rm bc}
=
\Pi_{13}K_X^2
\left(
c_{\rm bc}
+
w^2
+
\frac34w^4
\right)
K_X^2\Pi_{13}.
}
\end{equation}
This operator preserves the finite odd spatial slice and keeps compact winding
locked to the radial projection branch.

Define
\begin{equation}
\boxed{
q_1=\lambda_1-\lambda_0,
\qquad
q_3=\lambda_3-\lambda_0.
}
\end{equation}
The finite-rank radial datum is
\begin{equation}
\boxed{
\lambda_3
=
13.23388439508096.
}
\end{equation}
The parity-compatible leakage probability is
\begin{equation}
\boxed{
p_1
=
7.6528903366\times10^{-4},
\qquad
p_3=1-p_1.
}
\end{equation}
The boundary-resolved compact stiffness is
\begin{equation}
\boxed{
R_{A,\star}^{-2}
=
q_{\rm bc}
=
p_1q_1+(1-p_1)q_3.
}
\end{equation}
Using
\begin{equation}
\boxed{
\lambda_0=2.322863529580,
\qquad
\lambda_1=5.375830272676,
\qquad
\lambda_3=13.23388439508096,
}
\end{equation}
gives
\begin{equation}
\boxed{
q_{\rm bc}
=
10.905007182855176.
}
\end{equation}
The compact electromagnetic normalization is
\begin{equation}
\boxed{
\alpha_{\rm PR,bc}^{-1}
=
4\pi q_{\rm bc}
=
137.036361812007.
}
\end{equation}

The deterministic compact closure chain is
\begin{equation}
\boxed{
(T_A,T_X)
\rightarrow
(D_A,N_{\rm bc})
\rightarrow
c_{\rm bc}
\rightarrow
P_{\rm geom}^{\rm bc}
\rightarrow
p_1
\rightarrow
R_{A,\star}^{-2}
\rightarrow
\alpha_{\rm PR,bc}^{-1}.
}
\end{equation}
The electromagnetic normalization is an internal output of the compact boundary
map, not an experimental input.

\subsection{Magnetic Flux and Area-Law Diagnostic}

The compact phase flux normalization is
\begin{equation}
\boxed{
\Phi_\theta^{\rm PR}
=
\frac{\hbar}{e}
\frac{Z_A}{c_{\rm bc}}.
}
\end{equation}
With
\begin{equation}
Z_A=q_{\rm bc},
\end{equation}
this gives the numerical compact flux
\begin{equation}
\boxed{
\Phi_\theta^{\rm PR}
=
9.009597185\times10^{-15}\ {\rm Wb}.
}
\end{equation}
Equivalently,
\begin{equation}
\boxed{
\Phi_\theta^{\rm PR}
=
0.09009597185\ {\rm nG\,m^2}.
}
\end{equation}

For a realized projected compact magnetic domain, the field is determined by
the projected area:
\begin{equation}
\boxed{
B_{\rm geo}^{\rm PR}
=
\frac{
\Phi_\theta^{\rm PR}
}{
A_{\rm proj}^{(\theta)}
}.
}
\end{equation}
In nanogauss units,
\begin{equation}
\boxed{
B_{\rm geo}^{\rm PR}[{\rm nG}]
=
\frac{
0.09009597185
}{
A_{\rm proj}^{(\theta)}[{\rm m}^2]
}.
}
\end{equation}
This is an area law. The compact phase sector fixes the flux normalization, while
the realized projected compact area determines the magnetic amplitude. The theory
does not predict a universal fixed primordial magnetic-field amplitude.

\subsection{Closed Electromagnetic Projection Sequence}

The compact phase operator chain is
\begin{equation}
\boxed{
S^1_\theta
\rightarrow
O_\theta
=
-\frac{1}{R_A^2}
\frac{d^2}{d\theta^2}
\rightarrow
v_m(\theta)
=
\frac{e^{im\theta}}{\sqrt{2\pi R_A}}
\rightarrow
\lambda_m^{(\theta)}
=
\frac{m^2}{R_A^2}.
}
\end{equation}
The compact projection chain is
\begin{equation}
\boxed{
\Psi
\rightarrow
P_A[\Psi]
\rightarrow
\Phi_A(x)
=
A_\theta(x)e^{i\theta(x)}
\rightarrow
X_\mu
=
\partial_\mu\theta-q\mathcal A_\mu.
}
\end{equation}
The gauge curvature chain is
\begin{equation}
\boxed{
\mathcal A_\mu
\rightarrow
F_{\mu\nu}
=
\partial_\mu\mathcal A_\nu
-
\partial_\nu\mathcal A_\mu,
\qquad
F_{\mu\nu}'=F_{\mu\nu}.
}
\end{equation}
The compact kinetic chain is
\begin{equation}
\boxed{
G^{\theta\theta}=R_A^{-2}
\rightarrow
Z_A=R_A^{-2}
\rightarrow
\alpha_{\rm PR}^{-1}
=
4\pi Z_A.
}
\end{equation}
The boundary-resolved compact normalization chain is
\begin{equation}
\boxed{
(T_A,T_X)
\rightarrow
(D_A,N_{\rm bc})
\rightarrow
c_{\rm bc}
\rightarrow
P_{\rm geom}^{\rm bc}
\rightarrow
p_1
\rightarrow
R_{A,\star}^{-2}
\rightarrow
\alpha_{\rm PR,bc}^{-1}.
}
\end{equation}
The magnetic area-law chain is
\begin{equation}
\boxed{
\alpha_{\rm PR,bc}^{-1}
\rightarrow
Z_A
\rightarrow
\Phi_\theta^{\rm PR}
\rightarrow
B_{\rm geo}^{\rm PR}
=
\frac{\Phi_\theta^{\rm PR}}
     {A_{\rm proj}^{(\theta)}}.
}
\end{equation}

This completes the electromagnetic projection chain. The compact phase sector
generates the observable gauge connection, preserves local compact phase
covariance through \(X_\mu\), produces the gauge curvature \(F_{\mu\nu}\), fixes
the compact electromagnetic normalization through the boundary-resolved finite
rank map, and supplies the flux normalization used in the magnetic area-law
diagnostic.

\subsubsection{Compact-Phase Magnetic Area Law}

The compact phase closure fixes the electromagnetic winding stiffness and
therefore the compact flux normalization:
\[
\Phi_\theta^{\rm PR}
=
\frac{\hbar}{e}\frac{Z_A}{c_{bc}}.
\]
Numerically,
\[
\Phi_\theta^{\rm PR}
=
9.009597185\times10^{-15}\ {\rm Wb}.
\]
For a realized projected compact-phase magnetic domain with area
\(A_{\rm proj}^{(\theta)}\), the magnetic field is
\[
B_{\rm geo}^{\rm PR}
=
\frac{\Phi_\theta^{\rm PR}}
{A_{\rm proj}^{(\theta)}}.
\]
Equivalently,
\[
B_{\rm geo}^{\rm PR}[{\rm nG}]
=
\frac{0.09009597185}
{A_{\rm proj}^{(\theta)}[{\rm m^2}]}.
\]
The projected area is a physical property of the realized magnetic domain,
not an additional compact-sector parameter.

\subsubsection{Step 13: Electromagnetic Projection Summary}

The electromagnetic phase projection begins with the compact phase of the
master field:
\begin{equation}
\boxed{
P_A[\Psi](x)
=
\theta(x)
=
\arg
\left[
\langle U_A|\Psi_x\rangle_{\rm int}
\right].
}
\end{equation}

The observable compact phase gradient is
\begin{equation}
\boxed{
P_A[\Psi]_\mu
=
\partial_\mu\theta.
}
\end{equation}

Local compact phase coherence requires the gauge connection \(\mathcal A_\mu\):
\begin{equation}
\boxed{
X_\mu
=
\partial_\mu\theta
-
q\mathcal A_\mu.
}
\end{equation}

The electromagnetic field strength is
\begin{equation}
\boxed{
F_{\mu\nu}
=
\partial_\mu\mathcal A_\nu
-
\partial_\nu\mathcal A_\mu.
}
\end{equation}

The compact phase current is
\begin{equation}
\boxed{
J_\theta^\mu
=
A^2X^\mu.
}
\end{equation}

The current conservation law is
\begin{equation}
\boxed{
\nabla_\mu J_\theta^\mu
=
\nabla_\mu(A^2X^\mu)
=
0.
}
\end{equation}

The electromagnetic projection current is
\begin{equation}
\boxed{
J_{\rm EM}^\mu
=
-g_A^2qJ_\theta^\mu.
}
\end{equation}

The electromagnetic field equation is
\begin{equation}
\boxed{
\nabla_\mu F^{\mu\nu}
=
J_{\rm EM}^\nu.
}
\end{equation}

The compact gauge normalization is
\begin{equation}
\boxed{
g_A^{-2}
=
R_A^{-2}.
}
\end{equation}

The pre-boundary compact winding relation is
\begin{equation}
\boxed{
\alpha_{\rm PR}^{-1}
=
4\pi R_A^{-2}.
}
\end{equation}

The boundary-resolved fine-structure closure is
\begin{equation}
\boxed{
\alpha_{{\rm PR},bc}^{-1}
=
4\pi R_{A,\star}^{-2}
=
137.036361812007.
}
\end{equation}

Therefore, the complete electromagnetic projection chain is
\begin{equation}
\boxed{
\begin{aligned}
\Psi
&\rightarrow
P_A[\Psi]
\rightarrow
\theta(x)
\rightarrow
X_\mu
=
\partial_\mu\theta
-
q\mathcal A_\mu
\rightarrow
F_{\mu\nu}
\rightarrow
J^\mu_{\rm EM}
\\
&\rightarrow
g_A^{-2}=R_A^{-2}
\rightarrow
P_{\rm geom}^{\rm bc}
\rightarrow
R_{A,\star}^{-2}
\rightarrow
\alpha_{{\rm PR},bc}^{-1}.
\end{aligned}
}
\end{equation}

The intermediate compact stationarity reference is
\begin{equation}
\boxed{
\alpha_{\rm PR,geom}^{-1}
=
137.036361812010,
}
\end{equation}
while the boundary-resolved finite-rank closure gives
\begin{equation}
\boxed{
\alpha_{{\rm PR},bc}^{-1}
=
137.036361812007.
}
\end{equation}

The operational value used by the electromagnetic chain is the
boundary-resolved compact closure value.  No separate laboratory-reference
closure is introduced.  The electromagnetic sector is not introduced as an
independent gauge theory.  It is generated as the compact phase projection of
the master field, with its gauge connection, conserved phase current, Maxwell
equation, gauge kinetic normalization, and fine-structure normalization all
inherited from the internal compact phase geometry and its finite-rank
boundary-orientation closure.

\section{Projection Propagator and Spectral Self-Energy Chain}
\label{sec:projection_propagator_chain}

The projection propagator describes how residual internal spectral excitations
modify macroscopic propagation while preserving the massless low-energy pole.
This appendix records the derivation ledger for the subtracted self-energy
kernel, the corrected spin--2 propagator, the massless-pole residue, Newton
normalization, and the low-energy decoupling of projection-sector corrections.

The purpose of this chain is not to introduce an additional massive graviton.
The internal spectral response modifies the denominator through a gap-protected
self-energy. After the zero-momentum subtraction and Newton normalization, the
observable low-energy propagation remains massless, ghost-free, and luminal.

\subsection{Bare Massless Spin--2 Propagator}

In the transverse-traceless sector, the bare massless spin--2 propagator is
written as
\begin{equation}
\boxed{
D_{\mu\nu\rho\sigma}^{(0)}(k)
=
\frac{
P_{\mu\nu\rho\sigma}^{(2)}
}{
k^2+i\epsilon
}.
}
\end{equation}
Here \(P_{\mu\nu\rho\sigma}^{(2)}\) is the transverse-traceless spin--2
projector. The corresponding bare denominator is
\begin{equation}
\boxed{
P_0(k^2)=k^2.
}
\end{equation}

For compact notation, define
\begin{equation}
\boxed{
z=k^2.
}
\end{equation}

\subsection{Internal Spectral Density}

The macroscopic gravitational projection couples to excited internal modes above
the projection ground state. Let
\begin{equation}
\boxed{
\mu_a^2
=
\Lambda_a-\Lambda_0
}
\end{equation}
denote the shifted stiffness of the internal excitation labelled by \(a\), and let
\(g_a\) denote its coupling to the macroscopic projection channel. The internal
response density is
\begin{equation}
\boxed{
\rho_X(\mu^2)
=
\sum_{a\neq0}
|g_a|^2
\delta(\mu^2-\mu_a^2).
}
\end{equation}

The radial projection branch is gap protected:
\begin{equation}
\boxed{
\rho_X(\mu^2)=0,
\qquad
0\leq\mu^2<\mu_{\min}^2.
}
\end{equation}
The first radial gap is
\begin{equation}
\boxed{
\mu_{\min}^2
=
3.052966743096.
}
\end{equation}
Thus no arbitrarily soft internal excitation is available to modify the
low-energy massless pole.

\subsection{Bare Self-Energy Kernel}

The unrenormalized internal response contributes the bare self-energy
\begin{equation}
\boxed{
F_{\rm bare}(z)
=
\int_{\mu_{\min}^2}^{\infty}
d\mu^2\,
\frac{
\rho_X(\mu^2)
}{
z-\mu^2+i\epsilon
}.
}
\end{equation}
At zero external momentum,
\begin{equation}
F_{\rm bare}(0)
=
-\int_{\mu_{\min}^2}^{\infty}
d\mu^2\,
\frac{\rho_X(\mu^2)}{\mu^2}.
\end{equation}
For positive spectral weight, this is nonzero and negative:
\begin{equation}
\boxed{
\rho_X(\mu^2)\geq0
\quad\Longrightarrow\quad
F_{\rm bare}(0)<0.
}
\end{equation}
A direct use of \(F_{\rm bare}\) would shift the massless pole. The physical
kernel must therefore be subtracted at zero momentum.

\subsection{Massless-Pole Subtraction}

The physical self-energy is defined by subtracting the zero-momentum contribution:
\begin{equation}
\boxed{
F(z)
=
\int_{\mu_{\min}^2}^{\infty}
d\mu^2\,
\rho_X(\mu^2)
\left[
\frac{1}{z-\mu^2+i\epsilon}
+
\frac{1}{\mu^2}
\right].
}
\end{equation}
The bracket satisfies
\begin{equation}
\boxed{
\left[
\frac{1}{z-\mu^2}
+
\frac{1}{\mu^2}
\right]
=
\frac{z}{\mu^2(z-\mu^2)}.
}
\end{equation}
At \(z=0\),
\begin{equation}
\boxed{
F(0)=0.
}
\end{equation}
The subtraction therefore preserves the massless pole.

\subsection{Low-Energy Expansion}

For \(|z|\ll\mu_{\min}^2\), expand
\begin{equation}
\frac{1}{z-\mu^2}
=
-\frac{1}{\mu^2}
\left(
1+\frac{z}{\mu^2}
+
\frac{z^2}{\mu^4}
+
\cdots
\right).
\end{equation}
The subtracted bracket becomes
\begin{equation}
\frac{1}{z-\mu^2}
+
\frac{1}{\mu^2}
=
-\frac{z}{\mu^4}
-\frac{z^2}{\mu^6}
-\frac{z^3}{\mu^8}
-\cdots .
\end{equation}

Define the positive spectral moments
\begin{equation}
\boxed{
M_2
=
\int_{\mu_{\min}^2}^{\infty}
d\mu^2\,
\frac{\rho_X(\mu^2)}{\mu^4},
\qquad
M_3
=
\int_{\mu_{\min}^2}^{\infty}
d\mu^2\,
\frac{\rho_X(\mu^2)}{\mu^6},
\qquad
M_4
=
\int_{\mu_{\min}^2}^{\infty}
d\mu^2\,
\frac{\rho_X(\mu^2)}{\mu^8}.
}
\end{equation}
Then
\begin{equation}
\boxed{
F(z)
=
-zM_2
-z^2M_3
-z^3M_4
+O(z^4).
}
\end{equation}

\subsection{Corrected Propagator}

The corrected spin--2 propagator is
\begin{equation}
\boxed{
D_{\mu\nu\rho\sigma}^{(P)}(k)
=
\frac{
P_{\mu\nu\rho\sigma}^{(2)}
}{
k^2-F(k^2)+i\epsilon
}.
}
\end{equation}
Using \(z=k^2\), the corrected denominator is
\begin{equation}
P(z)
=
z-F(z).
\end{equation}
The low-energy expansion gives
\begin{equation}
P(z)
=
z
+
zM_2
+
z^2M_3
+
z^3M_4
+
O(z^4),
\end{equation}
or
\begin{equation}
\boxed{
P(z)
=
z
\left[
1+M_2+zM_3+z^2M_4+O(z^3)
\right].
}
\end{equation}

At the pole,
\begin{equation}
\boxed{
P(0)=0.
}
\end{equation}
The derivative at the pole is
\begin{equation}
\boxed{
P'(0)=1+M_2.
}
\end{equation}
Since \(M_2>0\),
\begin{equation}
\boxed{
P'(0)>0.
}
\end{equation}

\subsection{Massless-Pole Residue and Ghost Freedom}

The residue of the massless pole is
\begin{equation}
\boxed{
Z_{\rm pole}
=
\frac{1}{P'(0)}
=
\frac{1}{1+M_2}.
}
\end{equation}
Since \(M_2>0\),
\begin{equation}
\boxed{
0<Z_{\rm pole}<1.
}
\end{equation}
The massless pole remains present and has positive residue. The projection-sector
self-energy therefore does not introduce a ghost at the macroscopic massless
spin--2 pole.

\subsection{Newton Normalization}

The pole residue rescales the bare projection coupling \(G_P\). Matching the
observed Newton coupling gives
\begin{equation}
\boxed{
G_N
=
G_PZ_{\rm pole}
=
\frac{G_P}{1+M_2}.
}
\end{equation}
Equivalently,
\begin{equation}
\boxed{
G_P
=
(1+M_2)G_N.
}
\end{equation}
After this normalization, the constant pole-residue correction is absorbed into
the measured gravitational coupling. The remaining projection-sector correction
begins at higher order.

\subsection{Newton-Normalized Residual Kernel}

Define the Newton-normalized denominator by
\begin{equation}
\boxed{
z-F_E(z)
=
\frac{z-F(z)}{1+M_2}.
}
\end{equation}
Using the expansion of \(F(z)\),
\begin{equation}
z-F_E(z)
=
z
+
\frac{M_3}{1+M_2}z^2
+
\frac{M_4}{1+M_2}z^3
+
O(z^4).
\end{equation}
Hence
\begin{equation}
\boxed{
F_E(z)
=
-\frac{M_3}{1+M_2}z^2
-
\frac{M_4}{1+M_2}z^3
+
O(z^4).
}
\end{equation}
Therefore,
\begin{equation}
\boxed{
F_E(z)=O(z^2),
\qquad
\frac{F_E(z)}{z}\rightarrow0
\quad
(z\rightarrow0).
}
\end{equation}
In terms of \(k^2\),
\begin{equation}
\boxed{
F_E(k^2)=O(k^4),
\qquad
\frac{F_E(k^2)}{k^2}\rightarrow0
\quad
(k^2\rightarrow0).
}
\end{equation}

\subsection{Relative Projection Correction}

The relative projection-sector correction is
\begin{equation}
\boxed{
R_X(z)
=
\frac{F_E(z)}{z}.
}
\end{equation}
Because \(F_E(z)=O(z^2)\),
\begin{equation}
\boxed{
R_X(z)\rightarrow0
\qquad
(z\rightarrow0).
}
\end{equation}

The observable denominator correction may be written as
\begin{equation}
\boxed{
\Sigma_X(z)
=
\frac{F_E(z)}{z-F_E(z)}
=
\frac{R_X(z)}{1-R_X(z)}.
}
\end{equation}
If
\begin{equation}
\boxed{
|R_X(z)|\leq\delta_X<1,
}
\end{equation}
then
\begin{equation}
\boxed{
|\Sigma_X(z)|
\leq
\frac{\delta_X}{1-\delta_X}.
}
\end{equation}
In the weak projection regime,
\begin{equation}
\boxed{
\delta_X\ll1
\quad\Longrightarrow\quad
|\Sigma_X(z)|\ll1.
}
\end{equation}

\subsection{Euclidean Stability Check}

For spacelike Euclidean momentum, write
\begin{equation}
z=-Q^2,
\qquad
Q^2>0.
\end{equation}
The subtracted bracket becomes
\begin{equation}
\frac{1}{-Q^2-\mu^2}
+
\frac{1}{\mu^2}
=
\frac{Q^2}{\mu^2(Q^2+\mu^2)}.
\end{equation}
Thus
\begin{equation}
\boxed{
Q^2>0,
\qquad
\mu^2>0
\quad\Longrightarrow\quad
\frac{1}{-Q^2-\mu^2}
+
\frac{1}{\mu^2}
\geq0.
}
\end{equation}
For positive spectral density,
\begin{equation}
\boxed{
F(-Q^2)\geq0.
}
\end{equation}
The Euclidean response is bounded by the gap and does not create an unregulated
low-energy instability.

\subsection{Low-Energy Exterior Recovery}

The Newton-normalized propagator is
\begin{equation}
\boxed{
D_{\mu\nu\rho\sigma}^{(N)}(k)
=
\frac{
G_NP_{\mu\nu\rho\sigma}^{(2)}
}{
k^2-F_E(k^2)+i\epsilon
}.
}
\end{equation}
Since
\begin{equation}
\frac{F_E(k^2)}{k^2}\rightarrow0
\qquad
(k^2\rightarrow0),
\end{equation}
the low-energy propagator reduces to
\begin{equation}
\boxed{
D_{\mu\nu\rho\sigma}^{(N)}(k)
\rightarrow
\frac{
G_NP_{\mu\nu\rho\sigma}^{(2)}
}{
k^2+i\epsilon
}.
}
\end{equation}
The projection correction tensor satisfies
\begin{equation}
\boxed{
\Delta_{\mu\nu}^{(P)}
\rightarrow0
}
\end{equation}
in the low-energy exterior regime. The standard massless gravitational
propagation channel is recovered.

\subsection{Closed Propagator Stability Sequence}

The spectral-response chain is
\begin{equation}
\boxed{
\Psi
\rightarrow
P_X[\Psi]
\rightarrow
\rho_X(\mu^2)
=
\sum_{a\neq0}
|g_a|^2
\delta(\mu^2-\mu_a^2).
}
\end{equation}
The gap-protection chain is
\begin{equation}
\boxed{
\mu_{\min}^2>0
\rightarrow
\rho_X(\mu^2)=0
\quad
(0\leq\mu^2<\mu_{\min}^2).
}
\end{equation}
The subtraction chain is
\begin{equation}
\boxed{
F_{\rm bare}(z)
=
\int_{\mu_{\min}^2}^{\infty}
d\mu^2
\frac{\rho_X(\mu^2)}{z-\mu^2+i\epsilon}
\rightarrow
F(z)
=
\int_{\mu_{\min}^2}^{\infty}
d\mu^2\,
\rho_X(\mu^2)
\left[
\frac{1}{z-\mu^2+i\epsilon}
+
\frac{1}{\mu^2}
\right].
}
\end{equation}
The massless-pole chain is
\begin{equation}
\boxed{
F(0)=0
\rightarrow
P(0)=0
\rightarrow
Z_{\rm pole}
=
\frac{1}{1+M_2}
>
0.
}
\end{equation}
The Newton-normalization chain is
\begin{equation}
\boxed{
G_N
=
G_PZ_{\rm pole}
=
\frac{G_P}{1+M_2},
\qquad
F_E(z)=O(z^2),
\qquad
\frac{F_E(z)}{z}\rightarrow0.
}
\end{equation}
The exterior recovery chain is
\begin{equation}
\boxed{
z\ll\mu_{\min}^2
\rightarrow
\frac{F_E(z)}{z}\rightarrow0
\rightarrow
D_{\mu\nu\rho\sigma}^{(N)}(k)
\rightarrow
\frac{
G_NP_{\mu\nu\rho\sigma}^{(2)}
}{
k^2+i\epsilon
}
\rightarrow
\Delta_{\mu\nu}^{(P)}\rightarrow0.
}
\end{equation}

This completes the projection propagator chain. The internal spectral sector can
modify high-energy propagation through a gap-protected self-energy, but the
zero-momentum subtraction and Newton normalization preserve the massless
spin--2 pole, give a positive pole residue, remove ghost instabilities at low
energy, and recover ordinary gravitational propagation in the exterior
macroscopic limit.

\section{Gravitational-Wave Projection Chain}
\label{sec:gravitational_wave_projection_chain}

The gravitational-wave sector is the radiative projection of the determinant-normalized metric. This appendix records the chain by which the master-field stiffness projection produces Kerr exterior perturbations, inherits the Teukolsky operator in the exterior regime, and suppresses projection-sector residuals through the same gap-protected self-energy developed in the projection propagator chain.

The purpose of this section is not to introduce a new gravitational-wave mode hierarchy. In the exterior low-energy regime, Projection Relativity preserves the Kerr quasi-normal-mode hierarchy. The projection-sector contribution appears only as a gap-suppressed residual after Kerr subtraction.

\subsection{Metric Perturbation and Radiative Projection}

The gravitational projection produces the determinant-normalized effective metric
\begin{equation}
\boxed{
g_{\mu\nu}^{\rm eff}
=
\frac{Q_{\mu\nu}}{|\det Q|^{1/4}}.
}
\end{equation}
In a radiative exterior region, write the projected metric as a stationary background plus perturbation:
\begin{equation}
\boxed{
g_{\mu\nu}^{\rm eff}
=
\bar g_{\mu\nu}
+
h_{\mu\nu}.
}
\end{equation}
The gravitational-wave component is the transverse--traceless spin--2 projection of the perturbation:
\begin{equation}
\boxed{
h_{\mu\nu}^{\rm TT}
=
P_{\mu\nu}^{(2)\,\rho\sigma}
h_{\rho\sigma}.
}
\end{equation}
Here \(P^{(2)}\) denotes the spin--2 radiative projector. This isolates the propagating tensor response from gauge and non-radiative components.

\subsection{Exterior Kerr Inheritance}

Outside the finite projection core,
\begin{equation}
\boxed{
r>r_c,
}
\end{equation}
the exterior projection correction vanishes in the low-energy vacuum limit:
\begin{equation}
\boxed{
\Delta_{\mu\nu}^{(P)}
\rightarrow0.
}
\end{equation}
The exterior metric therefore reduces to Kerr:
\begin{equation}
\boxed{
\bar g_{\mu\nu}
=
g_{\mu\nu}^{\rm Kerr},
\qquad
r>r_c.
}
\end{equation}
Consequently, the leading perturbation operator is the ordinary spin-weighted Teukolsky operator,
\begin{equation}
\boxed{
O_{\rm Kerr}^{(s)}
=
O_{\rm Teuk}^{(s)}.
}
\end{equation}
The spin label \(s\) denotes the spin weight of the perturbation variable. Projection Relativity therefore inherits the Kerr perturbation hierarchy in the exterior regime rather than replacing it.

\subsection{Projection-Corrected Exterior Wave Operator}

The internal response sector adds a residual self-energy correction to the exterior wave operator. The corrected operator is
\begin{equation}
\boxed{
O_{\rm PR}^{(s)}
=
O_{\rm Teuk}^{(s)}
+
\widehat{\Sigma}_{X}^{\rm GW}.
}
\end{equation}
The gravitational-wave self-energy is the Newton-normalized projection residual evaluated on the Kerr operator:
\begin{equation}
\boxed{
\widehat{\Sigma}_{X}^{\rm GW}
=
-F_E\!\left(O_{\rm Teuk}^{(s)}\right).
}
\end{equation}
We find that
\begin{equation}
\boxed{
O_{\rm PR}^{(s)}
=
O_{\rm Teuk}^{(s)}
-
F_E\!\left(O_{\rm Teuk}^{(s)}\right).
}
\end{equation}

In a local wave patch, let
\begin{equation}
\boxed{
z
=
k^2.
}
\end{equation}
Then the corrected operator reduces locally to
\begin{equation}
\boxed{
O_{\rm PR}^{(s)}
\rightarrow
z-F_E(z).
}
\end{equation}

\subsection{Gap-Protected Low-Energy Recovery}

From the projection propagator chain, the internal spectral density is gap protected:
\begin{equation}
\boxed{
\rho_X(\mu^2)=0,
\qquad
0\leq\mu^2<\mu_{\min}^2.
}
\end{equation}
The radial spectral gap is
\begin{equation}
\boxed{
\mu_{\min}^2
=
3.052966743096.
}
\end{equation}
After zero-momentum subtraction and Newton normalization, the residual kernel satisfies
\begin{equation}
\boxed{
F_E(z)=O(z^2),
\qquad
\frac{F_E(z)}{z}\rightarrow0
\quad
(z\rightarrow0).
}
\end{equation}
Therefore,
\begin{equation}
\boxed{
O_{\rm PR}^{(s)}
\rightarrow
O_{\rm Teuk}^{(s)}
\qquad
(|z|\ll\mu_{\min}^2).
}
\end{equation}
This is the Kerr-recovery limit. The exterior low-energy operator remains the Kerr/Teukolsky operator, with projection-sector corrections suppressed by the internal spectral gap.

\subsection{Luminal Low-Energy Propagation}

The local dispersion denominator is
\begin{equation}
\boxed{
z-F_E(z).
}
\end{equation}
Since \(F_E(z)=O(z^2)\), the low-energy root remains
\begin{equation}
\boxed{
z=0.
}
\end{equation}
Thus the low-energy gravitational-wave dispersion relation is
\begin{equation}
\boxed{
\omega=c|k|.
}
\end{equation}
The projection-sector residual is not a massive-graviton contribution and does not introduce a low-energy gravitational-wave mass term.

\subsection{Response Function Expansion}

Let the sourced spin-weighted perturbation equation be
\begin{equation}
\boxed{
O_{\rm PR}^{(s)}\psi_{\rm PR}^{(s)}
=
S^{(s)}.
}
\end{equation}
The corresponding retarded Green operator is
\begin{equation}
\boxed{
G_{\rm PR}^{(s)}
=
\left[
O_{\rm Teuk}^{(s)}
+
\widehat{\Sigma}_{X}^{\rm GW}
\right]^{-1}.
}
\end{equation}
Define the Kerr Green operator by
\begin{equation}
\boxed{
G_{\rm Kerr}^{(s)}
=
\left[
O_{\rm Teuk}^{(s)}
\right]^{-1}.
}
\end{equation}
For a weak exterior projection correction,
\begin{equation}
\boxed{
\left\|
G_{\rm Kerr}^{(s)}
\widehat{\Sigma}_{X}^{\rm GW}
\right\|
\ll1,
}
\end{equation}
the Neumann expansion gives
\begin{equation}
\boxed{
G_{\rm PR}^{(s)}
=
G_{\rm Kerr}^{(s)}
-
G_{\rm Kerr}^{(s)}
\widehat{\Sigma}_{X}^{\rm GW}
G_{\rm Kerr}^{(s)}
+
O\!\left[
\left(
\widehat{\Sigma}_{X}^{\rm GW}
\right)^2
\right].
}
\end{equation}
Acting on the source,
\begin{equation}
\boxed{
\psi_{\rm PR}^{(s)}
=
\psi_{\rm Kerr}^{(s)}
+
\delta\psi_X^{(s)}
+
O\!\left[
\left(
\widehat{\Sigma}_{X}^{\rm GW}
\right)^2
\right],
}
\end{equation}
where
\begin{equation}
\boxed{
\delta\psi_X^{(s)}
=
-
G_{\rm Kerr}^{(s)}
\widehat{\Sigma}_{X}^{\rm GW}
\psi_{\rm Kerr}^{(s)}.
}
\end{equation}

\subsection{Observable Ringdown Decomposition}

For outgoing radiation, define the complex detector strain by
\begin{equation}
\boxed{
H(t)
=
h_+(t)-ih_\times(t).
}
\end{equation}
In frequency space, the radiative curvature perturbation induces a strain decomposition of the same form:
\begin{equation}
\boxed{
H_{\rm PR}(\omega)
=
H_{\rm Kerr}(\omega)
+
H_X(\omega).
}
\end{equation}
At the level of response amplitudes,
\begin{equation}
\boxed{
A_{\rm obs}(\omega)
=
A_{\rm Kerr}(\omega)
+
A_X(\omega).
}
\end{equation}
Here \(A_{\rm Kerr}\) is the ordinary Kerr ringdown response, including the Kerr mode hierarchy and any Kerr overtones used in the subtraction model. The term \(A_X\) is the residual projection-sector response. It is not a shifted Kerr fundamental mode and not a new low-energy pole.

The time-domain residual target is
\begin{equation}
\boxed{
h_X(t)
=
h_{\rm obs}(t)
-
h_{\rm Kerr}^{\rm full}(t),
}
\end{equation}
where \(h_{\rm Kerr}^{\rm full}\) denotes the full Kerr subtraction model used in the data analysis.

\subsection{Residual Suppression Bound}

The relative residual is controlled by the Newton-normalized self-energy ratio:
\begin{equation}
\boxed{
R_X(z)
=
\frac{F_E(z)}{z}.
}
\end{equation}
Since
\begin{equation}
\boxed{
R_X(z)\rightarrow0
\qquad
(z\rightarrow0),
}
\end{equation}
the projection-sector residual is suppressed in the low-energy exterior regime. The amplitude bound used in the main text is
\begin{equation}
\boxed{
\frac{\|A_X\|}{\|A_{\rm Kerr}\|}
\lesssim
C_{\rm Kerr}
\frac{\eta}{(1+M_2)(1-\eta)}
\epsilon_X^{\rm GW},
\qquad
\eta
=
\frac{|z|}{\mu_{\min}^2}.
}
\end{equation}
Where,
\begin{equation}
\boxed{
\eta\ll1
\quad\Longrightarrow\quad
\|A_X\|\ll\|A_{\rm Kerr}\|.
}
\end{equation}
This is the gravitational-wave consistency condition. The dominant exterior ringdown remains Kerr-like, and the projection-sector contribution is a suppressed residual channel.

\subsection{Finite-Core Restoration Interpretation}

The projection-sector residual is most naturally interpreted as a finite-core restoration response. The same radial operator that sets the finite-core spectral gap also sets the internal restoration spectrum:
\begin{equation}
\boxed{
O_X
=
-\frac{d^2}{dw^2}
+
1+w^2+\frac34w^4,
\qquad
\mathcal H_X
=
O_X-\lambda_0.
}
\end{equation}
The radial restoration gaps are
\begin{equation}
\boxed{
\mu_n^2
=
\lambda_n-\lambda_0.
}
\end{equation}
A general dissipative restoration operator can be written as
\begin{equation}
\boxed{
\widehat{\Gamma}_X
=
\left(
\mathcal C_X^{\rm eff}
\right)^{-1}
\mathcal H_X.
}
\end{equation}
The effective continuum/friction operator is determined by the imaginary part of the retarded internal self-energy:
\begin{equation}
\boxed{
\mathcal C_X^{\rm eff}(\omega_\ast)
=
-
\left.
\partial_\omega
{\rm Im}\,\Sigma_X^R(\omega)
\right|_{\omega=\omega_\ast}.
}
\end{equation}
In the first common-scale observational template,
\begin{equation}
\boxed{
\mathcal C_X^{\rm eff}=c_0I,
\qquad
\Gamma_n
=
\frac{\mu_n^2}{c_0}.
}
\end{equation}
This restoration interpretation does not replace Kerr. It gives a controlled way to parameterize a future high-signal-to-noise residual search after full Kerr subtraction.

\subsection{Current Observational Status}

The present manuscript does not claim a gravitational-wave detection of the projection-sector residual. Public-data screening with the fixed PR restoration template did not produce a statistically robust residual after comparison with generic drift and decay controls. The gravitational-wave channel is therefore a Kerr-consistency and future residual-stack channel in the present work.

\subsection{Closed Gravitational-Wave Projection Sequence}

The gravitational-wave projection chain is
\begin{equation}
\boxed{
\Psi
\rightarrow
\{U_{n,m},\Lambda_{n,m}\}
\rightarrow
P_g[\Psi]
\rightarrow
g_{\mu\nu}^{\rm eff}
\rightarrow
h_{\mu\nu}^{\rm TT}
\rightarrow
O_{\rm Teuk}^{(s)}
\rightarrow
O_{\rm PR}^{(s)}
=
O_{\rm Teuk}^{(s)}
+
\widehat{\Sigma}_{X}^{\rm GW}.
}
\end{equation}
The residual self-energy chain is
\begin{equation}
\boxed{
F_E(z)=O(z^2)
\rightarrow
\frac{F_E(z)}{z}\rightarrow0
\rightarrow
\widehat{\Sigma}_{X}^{\rm GW}
=
-F_E\!\left(O_{\rm Teuk}^{(s)}\right)
\rightarrow
O_{\rm PR}^{(s)}
\rightarrow
O_{\rm Teuk}^{(s)}.
}
\end{equation}
The response chain is
\begin{equation}
\boxed{
G_{\rm PR}^{(s)}
=
\left[
O_{\rm Teuk}^{(s)}
+
\widehat{\Sigma}_{X}^{\rm GW}
\right]^{-1}
\rightarrow
G_{\rm Kerr}^{(s)}
-
G_{\rm Kerr}^{(s)}
\widehat{\Sigma}_{X}^{\rm GW}
G_{\rm Kerr}^{(s)}
\rightarrow
A_{\rm obs}
=
A_{\rm Kerr}
+
A_X.
}
\end{equation}
The suppression chain is
\begin{equation}
\boxed{
\rho_X(\mu^2)=0
\quad
(0\leq\mu^2<\mu_{\min}^2)
\rightarrow
F_E(z)=O(z^2)
\rightarrow
\frac{\|A_X\|}{\|A_{\rm Kerr}\|}
\lesssim
C_{\rm Kerr}
\frac{\eta}{(1+M_2)(1-\eta)}
\epsilon_X^{\rm GW}
\rightarrow
\|A_X\|\ll\|A_{\rm Kerr}\|.
}
\end{equation}
The observational chain is
\begin{equation}
\boxed{
h_{\rm obs}
\rightarrow
h_{\rm Kerr}^{\rm full}
\rightarrow
h_X
=
h_{\rm obs}
-
h_{\rm Kerr}^{\rm full}
\rightarrow
\text{future high-SNR residual-stack test}.
}
\end{equation}

This completes the gravitational-wave projection chain. Projection Relativity inherits the Kerr exterior ringdown hierarchy, preserves luminal low-energy gravitational-wave propagation, and suppresses projection-sector residuals through the internal radial gap. The present work treats this sector as a consistency and constraint channel, not as a positive residual-detection channel.

\section{Unified Projection Energy Chain}
\label{sec:unified_projection_energy_chain}

The gravitational, displacement, compact electromagnetic, projection-response,
and homogeneous phase sectors are not energetically independent systems. They
are sector-resolved projections of the same master-field action. This appendix
collects the unified energy chain: the internal spectral expansion, the
four-dimensional effective Lagrangian, the sector decomposition, the projection
Hamiltonian, isolated energy conservation, and low-energy relativistic
kinematic recovery.

The purpose of this section is to show how the separately derived projection
channels recombine into one conserved master-field energy. Energy may be
redistributed among projection sectors, but the isolated total projection energy
is conserved.

\subsection{Master-Field Input}

The master field expands in the internal spectral basis as
\begin{equation}
\boxed{
\Psi(x,w,\theta)
=
\sum_{n,m}
c_{n,m}(x)U_{n,m}(w,\theta).
}
\end{equation}
The internal modes satisfy
\begin{equation}
\boxed{
O_{\rm int}U_{n,m}
=
\Lambda_{n,m}U_{n,m}.
}
\end{equation}
The internal spectral data provide the stiffness, compact winding, and response
scales used by the macroscopic projection sectors. The full projection action is
\begin{equation}
\boxed{
S_{\rm PR}
=
\int d^4x\,d\mu_{\rm int}\,
\sqrt{-g_{\rm eff}}\,
\mathcal L_{\rm PR}(x,\xi).
}
\end{equation}
The internal measure is
\begin{equation}
\boxed{
d\mu_{\rm int}
=
R_A\,dw\,d\theta .
}
\end{equation}

\subsection{Four-Dimensional Effective Lagrangian}

Macroscopic observers do not resolve the internal coordinates directly. The
observable four-dimensional effective Lagrangian density is obtained by
integrating over the internal manifold:
\begin{equation}
\boxed{
\mathcal L_{\rm PR}^{(4)}(x)
=
\int_{\mathcal M_{\rm int}}
d\mu_{\rm int}\,
\mathcal L_{\rm PR}(x,\xi).
}
\end{equation}

Projection onto the sector profiles gives the sector decomposition
\begin{equation}
\boxed{
\mathcal L_{\rm PR}^{(4)}
=
\mathcal L_g
+
\mathcal L_{\rm disp}
+
\mathcal L_A
+
\mathcal L_X
+
\mathcal L_\theta
+
\mathcal L_{\rm mix}.
}
\end{equation}
Here \(\mathcal L_g\) is the gravitational stiffness sector,
\(\mathcal L_{\rm disp}\) is the displacement-amplitude sector,
\(\mathcal L_A\) is the compact electromagnetic phase sector,
\(\mathcal L_X\) is the projection-response sector, \(\mathcal L_\theta\) is
the homogeneous compact-phase sector, and \(\mathcal L_{\rm mix}\) contains
allowed cross-sector coherence terms.

This decomposition does not introduce independent energies by hand. It is the
sector-resolved form of the same internally integrated master Lagrangian.

\subsection{Canonical Momentum and Hamiltonian Density}

Let \(q_I\) denote the macroscopic projected degrees of freedom appearing in
\(\mathcal L_{\rm PR}^{(4)}\). Their canonical momenta are
\begin{equation}
\boxed{
\Pi_I
=
\frac{\partial\mathcal L_{\rm PR}^{(4)}}{\partial\dot q_I}.
}
\end{equation}
For the master field notation, the canonical momentum conjugate to \(\Psi\) is
\begin{equation}
\boxed{
\Pi_\Psi
=
\frac{\partial\mathcal L_{\rm PR}^{(4)}}{\partial(\partial_0\Psi)}.
}
\end{equation}

The unified projection Hamiltonian density is the Legendre transform
\begin{equation}
\boxed{
\mathcal H_{\rm PR}
=
\sum_I
\Pi_I\dot q_I
-
\mathcal L_{\rm PR}^{(4)}.
}
\end{equation}
Equivalently, in master-field shorthand,
\begin{equation}
\boxed{
\mathcal H_{\rm PR}
=
\Pi_\Psi\partial_0\Psi
-
\mathcal L_{\rm PR}^{(4)}.
}
\end{equation}

The sector decomposition induces the Hamiltonian split
\begin{equation}
\boxed{
\mathcal H_{\rm PR}
=
\mathcal H_g
+
\mathcal H_{\rm disp}
+
\mathcal H_A
+
\mathcal H_X
+
\mathcal H_\theta
+
\mathcal H_{\rm mix}.
}
\end{equation}
The terms are not separate conserved universes. They are sector components of
one master Hamiltonian.

\subsection{Total Projection Energy}

The total projection energy on a spatial hypersurface \(\Sigma_t\) is
\begin{equation}
\boxed{
E_{\rm PR}(t)
=
\int_{\Sigma_t}
d^3x\,
\mathcal H_{\rm PR}.
}
\end{equation}
Using the sector Hamiltonian split,
\begin{equation}
\boxed{
E_{\rm PR}
=
E_g
+
E_{\rm disp}
+
E_A
+
E_X
+
E_\theta
+
E_{\rm mix}.
}
\end{equation}
Each term represents the energy carried by a projection sector or by coherent
cross-sector exchange. Only the total isolated energy is required to be
conserved.

\subsection{Noether Conservation}

When the internal geometric boundary data are fixed,
\begin{equation}
\boxed{
\partial_0R_A=0,
\qquad
\partial_0\mu_{\min}^2=0,
}
\end{equation}
and no external flux crosses the boundary of the system, the effective
four-dimensional Lagrangian has continuous time-translation symmetry. Noether
conservation gives
\begin{equation}
\boxed{
\frac{dE_{\rm PR}}{dt}
=
-\Phi_E,
}
\end{equation}
where \(\Phi_E\) is the net energy flux through the boundary of the chosen
spatial domain.

For an isolated projection system,
\begin{equation}
\boxed{
\Phi_E=0,
}
\end{equation}
so
\begin{equation}
\boxed{
\frac{dE_{\rm PR}}{dt}=0.
}
\end{equation}
Projection-sector exchange can redistribute energy among
\(\mathcal H_g,\mathcal H_{\rm disp},\mathcal H_A,\mathcal H_X,\mathcal H_\theta\),
and \(\mathcal H_{\rm mix}\), but the isolated master-field energy remains
constant.

\subsection{Low-Energy Sector Decoupling}

In the macroscopic exterior regime,
\begin{equation}
\boxed{
r>r_c,
\qquad
|k^2|\ll\mu_{\min}^2,
}
\end{equation}
the projection-response kernel obeys
\begin{equation}
\boxed{
\frac{F_E(k^2)}{k^2}\rightarrow0.
}
\end{equation}
The cross-sector coherence terms are then suppressed:
\begin{equation}
\boxed{
\mathcal L_{\rm mix}\rightarrow0
\qquad
(|k^2|\ll\mu_{\min}^2).
}
\end{equation}
The localized projection modes reduce to their canonical low-energy free-field
forms, with effective inertia supplied by the displacement vacuum:
\begin{equation}
\boxed{
\mathcal M_{\rm eff}
=
g_0\mathcal I_AA_0.
}
\end{equation}

\subsection{Projection-Mode Mass Scale}

Restoring physical units, a projected mode eigenvalue
\(\Lambda_{n,m}^{\rm phys}\) has dimensions
\begin{equation}
\boxed{
[\Lambda_{n,m}^{\rm phys}]
=
L^{-2}.
}
\end{equation}
The associated geometric mode mass scale is
\begin{equation}
\boxed{
\mathcal M_{n,m}^{\rm geom}
=
\frac{\hbar}{c}
\sqrt{\Lambda_{n,m}^{\rm phys}}.
}
\end{equation}
The physical projected mode equation is
\begin{equation}
\boxed{
\left(
\Box-\Lambda_{n,m}^{\rm phys}
\right)c_{n,m}(x)
=
0.
}
\end{equation}
Applying the standard quantum operators,
\begin{equation}
\boxed{
E\rightarrow i\hbar\partial_t,
\qquad
p\rightarrow -i\hbar\nabla,
}
\end{equation}
gives
\begin{equation}
\boxed{
E_{n,m}^2
=
|p|^2c^2
+
\hbar^2c^2\Lambda_{n,m}^{\rm phys}
=
|p|^2c^2
+
\left(
\mathcal M_{n,m}^{\rm geom}
\right)^2c^4.
}
\end{equation}

This identifies the geometric mass scale associated with a projected mode. The
matter-sector inertia used in macroscopic source terms is then obtained through
the displacement-overlap chain.

\subsection{Relativistic Kinematic Recovery}

For an isolated low-energy projection configuration, define the total
macroscopic momentum by
\begin{equation}
\boxed{
\mathbf P_{\rm PR}
=
\int_{\Sigma_t}
d^3x\,\boldsymbol{\mathcal P}_{\rm PR}.
}
\end{equation}
The corresponding invariant energy relation is
\begin{equation}
\boxed{
E_{\rm PR}^2
=
|\mathbf P_{\rm PR}|^2c^2
+
M_{\rm PR}^2c^4.
}
\end{equation}
In the rest frame,
\begin{equation}
\boxed{
\mathbf P_{\rm PR}=0,
\qquad
E_{\rm PR}^{(0)}
=
M_{\rm PR}c^2.
}
\end{equation}

Thus the unified projection Hamiltonian recovers the ordinary special-relativistic
kinematic invariant in the low-energy exterior limit.

\subsection{Closed Unified Energy Sequence}

The unified Lagrangian chain is
\begin{equation}
\boxed{
\Psi
\rightarrow
O_{\rm int}
\rightarrow
\{U_{n,m},\Lambda_{n,m}\}
\rightarrow
\mathcal L_{\rm PR}
\rightarrow
\mathcal L_{\rm PR}^{(4)}
=
\int_{\mathcal M_{\rm int}}
d\mu_{\rm int}\,
\mathcal L_{\rm PR}.
}
\end{equation}

The sector decomposition chain is
\begin{equation}
\boxed{
\mathcal L_{\rm PR}^{(4)}
\rightarrow
\mathcal L_g
+
\mathcal L_{\rm disp}
+
\mathcal L_A
+
\mathcal L_X
+
\mathcal L_\theta
+
\mathcal L_{\rm mix}.
}
\end{equation}

The Hamiltonian chain is
\begin{equation}
\boxed{
\mathcal L_{\rm PR}^{(4)}
\rightarrow
\Pi_I
=
\frac{\partial\mathcal L_{\rm PR}^{(4)}}{\partial\dot q_I}
\rightarrow
\mathcal H_{\rm PR}
=
\sum_I\Pi_I\dot q_I-\mathcal L_{\rm PR}^{(4)}
\rightarrow
E_{\rm PR}
=
\int_{\Sigma_t}
d^3x\,\mathcal H_{\rm PR}.
}
\end{equation}

The conservation chain is
\begin{equation}
\boxed{
\partial_0R_A=0,
\qquad
\partial_0\mu_{\min}^2=0,
\qquad
\Phi_E=0
\quad
\Longrightarrow
\quad
\frac{dE_{\rm PR}}{dt}=0.
}
\end{equation}

The low-energy kinematic chain is
\begin{equation}
\boxed{
|k^2|\ll\mu_{\min}^2
\rightarrow
\mathcal L_{\rm mix}\rightarrow0
\rightarrow
E_{\rm PR}^2
=
|\mathbf P_{\rm PR}|^2c^2
+
M_{\rm PR}^2c^4.
}
\end{equation}

The complete unified projection-energy sequence is
\begin{equation}
\boxed{
\Psi
\rightarrow
\{P_g,P_{\rm disp},P_A,P_X,P_\theta\}
\rightarrow
\mathcal L_{\rm PR}^{(4)}
\rightarrow
\mathcal H_{\rm PR}
\rightarrow
E_{\rm PR}
\rightarrow
\frac{dE_{\rm PR}}{dt}=0
\rightarrow
E_{\rm PR}^2
=
|\mathbf P_{\rm PR}|^2c^2
+
M_{\rm PR}^2c^4.
}
\end{equation}

This completes the unified projection energy chain. The gravitational stiffness,
displacement, compact electromagnetic phase, projection-response, and homogeneous
phase sectors are sector-resolved contributions to a single conserved
master-field energy. In the low-energy exterior regime, the sector coupling terms
decouple and the standard relativistic energy relation is recovered.

\section{Information-Preservation Chain}
\label{sec:information_preservation_chain}

The finite-core replacement of the classical singularity must preserve the
master-field information ledger while remaining compatible with the exterior
area-entropy bound. The purpose of this appendix section is to make that
distinction explicit. Projection Relativity does not claim that an exterior
observer has access to an unlimited hidden volume of states inside the core.
Instead, the full master-field state evolves unitarily in the projection Hilbert
space, while an exterior observer accesses only a boundary-projected,
coarse-grained image of that state.

The information-preservation chain therefore has two layers. The first layer is
the full internal spectral state, which remains norm-preserving and unitary. The
second layer is the exterior observational image, which is projected onto
horizon boundary data and is therefore area-bounded.

\subsection{Projection Hilbert Space and Master State}

The master field belongs to the projection Hilbert space
\begin{equation}
\boxed{
\mathcal H_P
=
L^2(\mathcal M_{\rm int},d\mu_{\rm int}).
}
\end{equation}
This is the space in which the full internal spectral ledger is defined. It is
larger than the data directly accessible to an exterior observer because it
contains the internal mode amplitudes of the master field.

Using the internal spectral basis, the state may be expanded as
\begin{equation}
\boxed{
|\Psi(t)\rangle
=
\sum_{n,m}
c_{n,m}(t)|U_{n,m}\rangle .
}
\end{equation}
The coefficients \(c_{n,m}(t)\) are the time-dependent entries of the internal
projection ledger. The basis states satisfy
\begin{equation}
\boxed{
\langle U_{j,k}|U_{n,m}\rangle_P
=
\delta_{jn}\delta_{km}.
}
\end{equation}
With this normalization, the full projection norm is
\begin{equation}
\boxed{
\|\Psi(t)\|_P^2
=
\langle\Psi(t)|\Psi(t)\rangle_P
=
\sum_{n,m}
|c_{n,m}(t)|^2.
}
\end{equation}
Preserving this norm is the basic information-preservation requirement for the
full master-field state.

\subsection{Hermitian Projection Hamiltonian}

For an isolated projection system, the master-field state evolves according to
\begin{equation}
\boxed{
i\hbar\frac{d}{dt}|\Psi(t)\rangle
=
\widehat H_P|\Psi(t)\rangle .
}
\end{equation}
The projection Hamiltonian is Hermitian:
\begin{equation}
\boxed{
\widehat H_P^\dagger
=
\widehat H_P.
}
\end{equation}
This condition is the mathematical statement that the full internal projection
ledger evolves without loss or gain of total probability.

The corresponding time-evolution operator is defined by
\begin{equation}
\boxed{
|\Psi(t)\rangle
=
\widehat U_P(t,t_0)|\Psi(t_0)\rangle .
}
\end{equation}
Hermiticity of \(\widehat H_P\) implies that \(\widehat U_P\) is unitary:
\begin{equation}
\boxed{
\widehat U_P^\dagger(t,t_0)
\widehat U_P(t,t_0)
=
I.
}
\end{equation}
This is the core information-preservation statement. The full state may be
redistributed among spectral modes, but it is not destroyed.

\subsection{Norm Conservation}

The norm of the master state is
\begin{equation}
\langle\Psi|\Psi\rangle_P.
\end{equation}
Taking its time derivative gives
\begin{align}
\frac{d}{dt}\langle\Psi|\Psi\rangle_P
&=
\langle\dot\Psi|\Psi\rangle_P
+
\langle\Psi|\dot\Psi\rangle_P .
\end{align}
Using the Schrödinger-type evolution equation,
\begin{equation}
|\dot\Psi\rangle
=
-\frac{i}{\hbar}
\widehat H_P|\Psi\rangle,
\qquad
\langle\dot\Psi|
=
\frac{i}{\hbar}
\langle\Psi|\widehat H_P ,
\end{equation}
the two terms cancel:
\begin{align}
\frac{d}{dt}\langle\Psi|\Psi\rangle_P
&=
\frac{i}{\hbar}
\langle\Psi|\widehat H_P|\Psi\rangle_P
-
\frac{i}{\hbar}
\langle\Psi|\widehat H_P|\Psi\rangle_P
\\
&=
0.
\end{align}
Therefore,
\begin{equation}
\boxed{
\frac{d}{dt}\|\Psi(t)\|_P^2=0.
}
\end{equation}
Equivalently,
\begin{equation}
\boxed{
\sum_{n,m}|c_{n,m}(t)|^2
=
\sum_{n,m}|c_{n,m}(t_0)|^2.
}
\end{equation}
The spectral coefficients may move between modes, but the total ledger norm is
fixed.

\subsection{Density-Matrix Evolution}

The same conclusion can be stated in density-matrix language. For a pure
master-field state,
\begin{equation}
\boxed{
\rho(t)
=
|\Psi(t)\rangle\langle\Psi(t)|.
}
\end{equation}
Under unitary projection evolution,
\begin{equation}
\boxed{
\rho(t)
=
\widehat U_P(t,t_0)
\rho(t_0)
\widehat U_P^\dagger(t,t_0).
}
\end{equation}
The trace of the density matrix is preserved because cyclicity of the trace and
unitarity give
\begin{align}
{\rm Tr}\,\rho(t)
&=
{\rm Tr}
\left[
\widehat U_P\rho(t_0)\widehat U_P^\dagger
\right]
\\
&=
{\rm Tr}
\left[
\rho(t_0)\widehat U_P^\dagger\widehat U_P
\right]
\\
&=
{\rm Tr}\,\rho(t_0).
\end{align}
Hence
\begin{equation}
\boxed{
{\rm Tr}\,\rho(t)
=
{\rm Tr}\,\rho(t_0).
}
\end{equation}

Purity is also preserved under the full master-field evolution:
\begin{align}
{\rm Tr}\,\rho^2(t)
&=
{\rm Tr}
\left[
\widehat U_P\rho(t_0)\widehat U_P^\dagger
\widehat U_P\rho(t_0)\widehat U_P^\dagger
\right]
\\
&=
{\rm Tr}
\left[
\widehat U_P\rho^2(t_0)\widehat U_P^\dagger
\right]
\\
&=
{\rm Tr}\,\rho^2(t_0).
\end{align}
Therefore,
\begin{equation}
\boxed{
{\rm Tr}\,\rho^2(t)
=
{\rm Tr}\,\rho^2(t_0).
}
\end{equation}
A pure full state remains pure under the complete projection evolution. Any
apparent loss of purity must therefore arise from projection, coarse graining, or
restriction to an observer-accessible subsystem.

\subsection{Projection-Sector Coarse Graining}

An observer rarely has access to the full internal projection ledger. Instead,
the observer interacts with a restricted set of projection sectors. Let
\(\{\Pi_A\}\) be a complete orthogonal set of sector projectors:
\begin{equation}
\boxed{
\Pi_A\Pi_B
=
\delta_{AB}\Pi_A,
\qquad
\sum_A\Pi_A=I.
}
\end{equation}
The sector-resolved coarse-grained state is
\begin{equation}
\boxed{
\rho_{\rm cg}
=
\sum_A
\Pi_A\rho\Pi_A.
}
\end{equation}
This operation removes coherence between distinct projection sectors while
retaining the total probability assigned to the sectors.

The discarded off-sector coherence is
\begin{equation}
\boxed{
\rho_{\rm coh}
=
\rho-\rho_{\rm cg}
=
\sum_{A\neq B}
\Pi_A\rho\Pi_B.
}
\end{equation}
The trace of the coarse-grained state remains equal to the trace of the full
state:
\begin{align}
{\rm Tr}\,\rho_{\rm cg}
&=
{\rm Tr}
\left[
\sum_A
\Pi_A\rho\Pi_A
\right]
\\
&=
\sum_A
{\rm Tr}
\left[
\Pi_A\rho\Pi_A
\right]
\\
&=
\sum_A
{\rm Tr}
\left[
\rho\Pi_A^2
\right]
\\
&=
\sum_A
{\rm Tr}
\left[
\rho\Pi_A
\right]
\\
&=
{\rm Tr}
\left[
\rho
\sum_A\Pi_A
\right]
\\
&=
{\rm Tr}\,\rho.
\end{align}
Thus,
\begin{equation}
\boxed{
{\rm Tr}\,\rho_{\rm cg}
=
{\rm Tr}\,\rho.
}
\end{equation}
The observer has not lost total probability. What is lost is access to the
off-sector coherence.

For a positive density matrix, projection-sector coarse graining cannot increase
purity:
\begin{equation}
\boxed{
{\rm Tr}\,\rho_{\rm cg}^2
\leq
{\rm Tr}\,\rho^2.
}
\end{equation}
A strict inequality,
\begin{equation}
\boxed{
{\rm Tr}\,\rho_{\rm cg}^2
<
{\rm Tr}\,\rho^2,
}
\end{equation}
means that the restricted observer sees decoherence even though the full
master-field evolution remains unitary. This is the distinction needed for the
black-hole information problem: exterior decoherence does not imply fundamental
information loss.

\subsection{Finite-Core Spectral Redistribution}

The finite core appears when the classical collapse channel reaches the radial
projection ceiling:
\begin{equation}
\boxed{
\mu_{\min}^2>0
\quad\Longrightarrow\quad
R_{\max}<\infty
\quad\Longrightarrow\quad
r_c>0.
}
\end{equation}
The classical endpoint \(r=0\) is not part of the PR geometric domain. It is
replaced by a finite spectral core:
\begin{equation}
\boxed{
r\rightarrow r_c.
}
\end{equation}
At this point, the collapsing exterior description is replaced by spectral
redistribution inside the projection Hilbert space:
\begin{equation}
\boxed{
\{c_{n,m}(t_0)\}
\rightarrow
\{c_{a}^{(c)}(t_c),c_{b}^{({\rm ext})}(t_c)\}.
}
\end{equation}
The superscript \((c)\) labels finite-core spectral coefficients, while
\(({\rm ext})\) labels coefficients still represented in the exterior
projection channel.

The total spectral norm is preserved:
\begin{equation}
\boxed{
\sum_{n,m}
|c_{n,m}(t_0)|^2
=
\sum_a
|c_a^{(c)}(t_c)|^2
+
\sum_b
|c_b^{({\rm ext})}(t_c)|^2.
}
\end{equation}
Collapse is therefore not modeled as destruction of coefficients. It is modeled
as redistribution of the same master-field ledger between finite-core and
exterior projection sectors.

The finite-core spectral ledger is denoted
\begin{equation}
\boxed{
\mathcal B_c
=
\left\{
c_a^{(c)},\mu_a^2,\Pi_a^{(c)}
\right\}.
}
\end{equation}
This ledger records the bounded internal state after the curvature ceiling is
reached. Since the evolution operator is unitary, the initial state can be
reconstructed from the full final master-field state:
\begin{equation}
\boxed{
|\Psi(t_0)\rangle
=
\widehat U_P^\dagger(t_c,t_0)
|\Psi(t_c)\rangle.
}
\end{equation}
In density-matrix form,
\begin{equation}
\boxed{
\rho(t_0)
=
\widehat U_P^\dagger(t_c,t_0)
\rho(t_c)
\widehat U_P(t_c,t_0).
}
\end{equation}
The finite core therefore stores information in the internal spectral ledger
rather than erasing it at a classical singular endpoint.

\subsection{Core Ledger and Energy Split}

The finite core carries a bounded share of the projection energy. Its energy is
\begin{equation}
\boxed{
E_{\rm PR}^{(c)}
=
\int_{r\leq r_c}
d^3x\,
\mathcal H_{\rm PR}.
}
\end{equation}
The exterior projection channel carries
\begin{equation}
\boxed{
E_{\rm PR}^{({\rm ext})}
=
\int_{r>r_c}
d^3x\,
\mathcal H_{\rm PR}.
}
\end{equation}
The total isolated projection energy decomposes as
\begin{equation}
\boxed{
E_{\rm PR}
=
E_{\rm PR}^{(c)}
+
E_{\rm PR}^{({\rm ext})}.
}
\end{equation}
For an isolated system,
\begin{equation}
\boxed{
\frac{dE_{\rm PR}}{dt}=0.
}
\end{equation}
This gives
\begin{equation}
\boxed{
\frac{d}{dt}
\left[
E_{\rm PR}^{(c)}
+
E_{\rm PR}^{({\rm ext})}
\right]
=
0.
}
\end{equation}
Energy may move between the exterior projection and the finite-core ledger, but
the isolated total is conserved. The finite core stores energy in bounded
spectral form; it does not require divergent curvature or divergent density.

\subsection{Exterior Horizon Projection}

The finite core is not an additional exterior entropy volume. The full spectral
ledger belongs to the master-field Hilbert space, while an exterior observer
sees only a boundary-projected image of that ledger.

The horizon projection channel is
\begin{equation}
\boxed{
\mathcal E_H:
\mathcal B_c
\longrightarrow
\mathcal H_\partial.
}
\end{equation}
Here \(\mathcal B_c\) is the finite-core spectral ledger and
\(\mathcal H_\partial\) is the horizon boundary data. This map expresses that
the exterior image of the core is encoded on the boundary, not as an externally
accessible bulk volume.

The exterior entropy is bounded by the horizon area:
\begin{equation}
\boxed{
S_{\rm ext}
\leq
\frac{k_BA_H}{4\ell_P^2}.
}
\end{equation}
This is a statement about the coarse-grained exterior channel. It does not
replace the full internal unitary ledger. The area bound constrains what the
exterior observer can access, while the full state remains encoded in
\(\mathcal H_P\).

\subsection{Information Preservation and Exterior Coarse Graining}

Let \(\rho_c\) be the finite-core density matrix. The exterior boundary state is
the projected image
\begin{equation}
\boxed{
\rho_{\partial}
=
\mathcal E_H(\rho_c).
}
\end{equation}
The exterior entropy is
\begin{equation}
\boxed{
S_{\rm ext}
=
-k_B{\rm Tr}
\left(
\rho_{\partial}\ln\rho_{\partial}
\right).
}
\end{equation}
This entropy is the entropy of the boundary-projected image, not the entropy of
the complete master-field state. A full state may remain pure while its exterior
projection appears mixed.

The full master-field evolution preserves
\begin{equation}
\boxed{
\|\Psi(t)\|_P^2,
\qquad
{\rm Tr}\,\rho(t),
\qquad
{\rm Tr}\,\rho^2(t).
}
\end{equation}
An exterior observer generally accesses only
\begin{equation}
\boxed{
\rho_{\partial}
\neq
\rho.
}
\end{equation}
The difference between these two states is the difference between full
projection information and exterior coarse-grained information.

\subsection{Closed Information-Preservation Sequence}

The unitary evolution chain is
\begin{equation}
\boxed{
\widehat H_P^\dagger=\widehat H_P
\rightarrow
\widehat U_P^\dagger\widehat U_P=I
\rightarrow
\frac{d}{dt}\|\Psi\|_P^2=0
\rightarrow
\sum_{n,m}|c_{n,m}(t)|^2
=
{\rm constant}.
}
\end{equation}
This chain states that the master-field ledger is preserved at the level of the
full internal spectral state.

The density-matrix chain is
\begin{equation}
\boxed{
\rho(t)
=
\widehat U_P\rho(t_0)\widehat U_P^\dagger
\rightarrow
{\rm Tr}\,\rho(t)
=
{\rm Tr}\,\rho(t_0)
\rightarrow
{\rm Tr}\,\rho^2(t)
=
{\rm Tr}\,\rho^2(t_0).
}
\end{equation}
This is the density-matrix version of the same unitary preservation statement.

The projection-sector coarse-graining chain is
\begin{equation}
\boxed{
\rho
\rightarrow
\rho_{\rm cg}
=
\sum_A
\Pi_A\rho\Pi_A
\rightarrow
{\rm Tr}\,\rho_{\rm cg}
=
{\rm Tr}\,\rho,
\qquad
{\rm Tr}\,\rho_{\rm cg}^2
\leq
{\rm Tr}\,\rho^2.
}
\end{equation}
The trace is preserved under coarse graining, while purity can decrease because
sector coherence is no longer visible. The finite-core redistribution chain is
\begin{equation}
\boxed{
\begin{aligned}
\mu_{\min}^2 > 0 &\implies R_{\max} < \infty \implies r_c > 0 \\
&\implies \text{spectral redistribution as } r \to r_c \\
&\implies \|\Psi(t_c)\|_P^2 = \|\Psi(t_0)\|_P^2
\end{aligned}
}
\end{equation}
This replaces singular information loss with finite spectral redistribution. The exterior entropy chain is
\begin{equation}
\boxed{
\mathcal B_c
\rightarrow
\mathcal E_H(\mathcal B_c)
=
\mathcal H_\partial
\rightarrow
S_{\rm ext}
\leq
\frac{k_BA_H}{4\ell_P^2}.
}
\end{equation}
The exterior observer receives an area-bounded boundary image of the finite-core
ledger. The complete information-preservation chain is
\begin{equation}
\boxed{
\Psi
\rightarrow
\widehat H_P
\rightarrow
\widehat U_P
\rightarrow
\rho(t)
=
\widehat U_P\rho(t_0)\widehat U_P^\dagger
\rightarrow
{\rm Tr}\,\rho
=
{\rm constant}
\rightarrow
{\rm Tr}\,\rho^2
=
{\rm constant}
\rightarrow
\mathcal B_c
\rightarrow
\mathcal H_\partial.
}
\end{equation}
Including finite-core saturation, the final chain is
\begin{equation}
\boxed{
\mu_{\min}^2>0
\rightarrow
R_{\max}<\infty
\rightarrow
r_c>0
\rightarrow
\text{collapse as finite spectral redistribution}
\rightarrow
\|\Psi(t)\|_P^2
=
\|\Psi(t_0)\|_P^2.
}
\end{equation}

This completes the information-preservation chain. Projection Relativity
replaces the classical singular endpoint with a finite-core spectral state whose
full master-field evolution remains unitary. The exterior observer sees only an
area-bounded, coarse-grained boundary image of that unitary ledger.

\end{document}
